\documentclass[10pt,twoside]{siamart1116}

\usepackage[english]{babel}
\usepackage{fancyhdr,amsfonts,amssymb}
\usepackage{array}
\usepackage{xcolor}

\setlength{\textheight}{210mm}
\setlength{\textwidth}{165mm}
\topmargin = -10mm
\renewcommand{\baselinestretch}{1.1}
\setlength{\parskip}{.1in}

\newcommand{\CA}{Eric Darve}
\newcommand{\Names}{Eric Darve}
\newcommand{\Title}{A Unified Theory for the Existence of LU and $LDL^T$ Factorizations without Pivoting}

\renewcommand{\theequation}{\thesection.\arabic{equation}}

\usepackage{aliascnt}
% \siamthmalias{<env>}{<Name>}{<\newtheorem or \renewtheorem>}
\newcommand{\siamthmalias}[3]{%
  \newaliascnt{#1}{theorem}%
  \theoremstyle{plain}%
  \theoremheaderfont{\normalfont\sc}%
  \theorembodyfont{\normalfont\itshape}%
  \theoremseparator{.}%
  \theoremsymbol{}%
  #3{#1}{#2}%
  \expandafter\def\csname the#1\endcsname{\thetheorem}%
}
\siamthmalias{property}{Property}{\newtheorem}
\siamthmalias{lemma}{Lemma}{\renewtheorem}
\siamthmalias{corollary}{Corollary}{\renewtheorem}
\siamthmalias{proposition}{Proposition}{\renewtheorem}
\siamthmalias{definition}{Definition}{\renewtheorem}
\siamthmalias{remark}{Remark}{\newtheorem}
\siamthmalias{example}{Example}{\newtheorem}

\usepackage{booktabs,listings}
\usepackage{csquotes}
\MakeOuterQuote{"}

\lstdefinestyle{PythonStyle}{
  language=Python,
  basicstyle=\ttfamily\fontsize{7pt}{8pt}\selectfont,
  keywordstyle=\color{blue}\bfseries,
  stringstyle=\color{red},
  commentstyle=\color{green}\itshape,
  showstringspaces=false,
  numbers=left,
  numberstyle=\tiny\color{gray},
  breaklines=true,
  captionpos=b,
  escapeinside={(*@}{@*)},
}

\DeclareMathOperator{\rk}{rank}
\DeclareMathOperator{\nl}{null}
\DeclareMathOperator{\spn}{span}
\DeclareMathOperator{\range}{range}
\DeclareMathOperator{\sgn}{sgn}

\newcommand{\by}{\times}
\newcommand{\M}{\mathcal{M}}
\newcommand{\F}{F}

% Treat ":" as an ordinary atom in math mode so that slicing such as
% A[1:k,1:k] is set tightly.
\mathcode`\:="003A

\Crefname{property}{Property}{Properties}
\crefname{property}{property}{properties}
\Crefname{corollary}{Corollary}{Corollaries}
\crefname{corollary}{corollary}{corollaries}
\Crefname{lemma}{Lemma}{Lemmas}
\crefname{lemma}{lemma}{lemmas}
\Crefname{proposition}{Proposition}{Propositions}
\crefname{proposition}{proposition}{propositions}
\Crefname{theorem}{Theorem}{Theorems}
\crefname{theorem}{theorem}{theorems}
\Crefname{definition}{Definition}{Definitions}
\crefname{definition}{definition}{definitions}
\Crefname{remark}{Remark}{Remarks}
\crefname{remark}{remark}{remarks}
\Crefname{example}{Example}{Examples}
\crefname{example}{example}{examples}
\Crefname{lstlisting}{Listing}{Listings}
\crefname{lstlisting}{listing}{listings}

\numberwithin{theorem}{section}

\begin{document}

\setcounter{page}{1}
\thispagestyle{empty}

\title{\Title\thanks{Corresponding Author: \CA}}

\author{
  Eric Darve\thanks{Institute for Computational and Mathematical Engineering,
    Mechanical Engineering Department,
    Stanford University,
    Stanford, CA 94305, USA
  (darve@stanford.edu).}
}

\maketitle

\begin{abstract}
  We give a unified treatment of two existence questions for triangular
  factorizations without pivoting: when does a square matrix $A$ with entries
  in an arbitrary field $\F$ admit a factorization $A = LU$ with $L$ lower
  triangular and $U$ upper triangular, and, when $A$ is symmetric, a
  factorization $A = LDL^T$ with $L$ lower triangular and $D$ diagonal?  No
  permutations are allowed, $A$ may be singular, and zeros are allowed on the
  diagonals of all factors; this distinguishes the problem from the classical
  theory, which assumes that all leading principal minors are nonzero.  Both
  questions are answered by nullity inequalities comparing each leading
  principal submatrix with the corresponding leading block of columns and
  rows.  The unsymmetric criterion recovers the theorem of Okunev and Johnson;
  the symmetric criterion appears to be new.  Both sufficiency proofs follow a
  single strategy: reduce $A$ by an invertible triangular equivalence (a
  congruence in the symmetric case) to a \emph{matching normal form} --- a
  matrix whose nonzero pattern is a matching, unique for a given $A$ and
  closely related to the Bruhat decomposition --- observe that the criterion
  is invariant under such reductions and becomes a simple combinatorial prefix
  condition on the matching, and then factor the normal form explicitly by a
  greedy slot assignment (unsymmetric case) or by threading telescoping chains
  through the isolated indices (symmetric case).  The proofs are constructive
  and yield an $O(n^3)$ factorization algorithm.  We validate all results
  with exact arithmetic over $\mathbb{Q}$, $GF(2)$, and $GF(3)$, including
  exhaustive comparisons with brute-force enumeration of all triangular
  factorizations for small dimensions.
\end{abstract}

\begin{keywords}
  LU factorization, $LDL^T$ factorization, singular matrices, rank-deficient
  matrices, pivoting, nullity, Bruhat decomposition, matching.
\end{keywords}
\begin{AMS}
  15A23, 15A03, 15A21, 65F05.
\end{AMS}

\section{Introduction}\label{sec:intro}

Let $\F$ be an arbitrary field and let $A \in \F^{n \by n}$.  It is classical
that when $A$ is nonsingular, a factorization $A = LU$ with $L$ lower
triangular and $U$ upper triangular exists if and only if every leading
principal submatrix $A[1:k,1:k]$ is nonsingular~\cite{higham_gaussian_2011,golub_matrix_2013}.
The standard remedy when this condition fails is pivoting: for any $A$ there
are permutations $P$ (and $Q$) such that $PA = LU$ or $PAQ = LU$.  Pivoting,
however, changes the problem.  When the row and column indices carry meaning
--- temporal order in autoregressive models, causal order in structural
equation models, or a fixed elimination order imposed by an application ---
the factorization must respect the original indexing, and the correct
question becomes: \emph{for which matrices $A$ does $A = LU$ hold with no
permutations at all?}  We refer to~\cite{darve_lu_2026} for an extended
discussion of this motivation.

Two features make the unpivoted question fundamentally different from the
classical one.  First, $A$ may be singular, or may have singular leading
principal blocks, and yet admit a perfectly good factorization.  Second, the
factors must then be allowed to carry \emph{zeros on their diagonals}: we
require only that $L$ be lower triangular and $U$ upper triangular, not that
they be invertible.  The classical minor criterion says nothing in this
regime.  For example,
\begin{equation}\label{eq:swap-vs-exchange}
  A_1 =
  \begin{pmatrix}
    0 & 1\\
    1 & 0
  \end{pmatrix},
  \qquad
  A_2 =
  \begin{pmatrix}
    0 & 0 & 0\\
    0 & 0 & 1\\
    0 & 1 & 0
  \end{pmatrix}
\end{equation}
are both symmetric with singular leading blocks.  The swap matrix $A_1$ has
no factorization $A_1 = LU$ (from $a_{11} = l_{11}u_{11} = 0$, either the
first row or the first column of $LU$ must vanish).  Yet the larger, more
degenerate matrix $A_2$ factors:
\begin{equation}\label{eq:exchange-factorizations}
  A_2 =
  \underbrace{\begin{pmatrix} 0 & 0 & 0\\ 0 & 1 & 0\\ 1 & 0 & 0 \end{pmatrix}
  \begin{pmatrix} 0 & 1 & 0\\ 0 & 0 & 1\\ 0 & 0 & 0 \end{pmatrix}}_{LU}
  =
  \underbrace{\begin{pmatrix} 0 & 0 & 0\\ 1 & 1 & 0\\ 0 & 1 & 1 \end{pmatrix}
  \begin{pmatrix} -1 & & \\ & 1 & \\ & & -1 \end{pmatrix}
  \begin{pmatrix} 0 & 1 & 0\\ 0 & 1 & 1\\ 0 & 0 & 1 \end{pmatrix}}_{LDL^T} .
\end{equation}
Both factorizations essentially require zero diagonal entries: the leading
zero index of $A_2$ acts as a resource that later off-diagonal structure can
be routed through.  Understanding exactly when such routing is possible is
the subject of this paper.

We answer the question for both the unsymmetric and the symmetric form of
the problem, over an arbitrary field, in a unified way.  Write
$\nl(X) = (\#\text{cols of } X) - \rk(X)$ for the (right) nullity.  Our two
main results are the following (\Cref{thm:lu,thm:ldlt}):
\begin{itemize}
  \item $A \in \F^{n \by n}$ admits a factorization $A = LU$ with $L$ lower
    triangular and $U$ upper triangular, zeros allowed on both diagonals, if
    and only if for every $k = 1, \ldots, n$,
    \[
      \nl(A[1:k,1:k]) \le \nl(A[:,1:k]) + \nl({A[1:k,:]}^T) .
    \]
  \item A symmetric $A \in \F^{n \by n}$ admits a factorization $A = LDL^T$
    with $L$ lower triangular (zeros allowed on the diagonal) and $D$
    diagonal if and only if for every $k = 1, \ldots, n$,
    \[
      \nl(A[1:k,1:k]) \le 2 \, \nl(A[:,1:k]) .
    \]
\end{itemize}
The first criterion is equivalent to the rank inequality of Okunev and
Johnson~\cite{okunev_johnson_2005,johnson_okunev_2025}, restated in the
nullity form introduced in~\cite{darve_lu_2026}.  The second criterion, and
the fact that for symmetric matrices the two criteria coincide --- so that a
symmetric matrix admits an unpivoted $LU$ factorization if and only if it
admits an unpivoted $LDL^T$ factorization --- appear to be new.

\paragraph{The unified method}
Both sufficiency proofs follow the same three-step strategy, which is the
main methodological contribution of this paper.
\begin{enumerate}
  \item \emph{Reduce.}  Every $A \in \F^{n \by n}$ can be written
    $A = M \Pi N$ with $M$ invertible lower triangular, $N$ invertible upper
    triangular, and $\Pi$ a \emph{matching matrix}: a matrix with at most one
    nonzero entry in every row and every column
    (\Cref{lem:reduction-lu}).  Every symmetric $A$ can be written
    $A = M P M^T$ with $M$ invertible lower triangular and $P$ a
    \emph{symmetric matching matrix}, whose nonzero blocks are $1 \by 1$
    diagonal entries and $2 \by 2$ blocks
    $\bigl(\begin{smallmatrix} 0 & \alpha\\ \alpha & \beta
    \end{smallmatrix}\bigr)$ (\Cref{lem:reduction-sym}); over $\mathbb{R}$
    the normal form can be scaled to a \emph{signed matching matrix} with
    entries in $\{0, \pm 1\}$ (\Cref{rem:signed}), recovering the
    unpivoted symmetric elimination of Van Dooren and
    Strang~\cite{vandooren_strang_2014}.  The nonzero pattern of the
    normal form is uniquely determined by $A$ through its rank profile
    (\Cref{prop:canonical}); for invertible $A$ this is the Bruhat
    decomposition, and in general it is the rank profile
    matrix~\cite{dumas_pernet_sultan_2017} and its congruence analogue
    (cf.~\cite{szechtman_2007}).
  \item \emph{Transfer.}  All the nullities appearing in the criteria are
    invariant under $A \mapsto MAN$ with $M, N$ invertible triangular
    (\Cref{lem:transfer}).  Hence $A$ satisfies a criterion if and only if
    its normal form does, and on the normal form the criterion collapses to
    a purely combinatorial \emph{prefix condition} on the matching
    (\Cref{lem:comb-lu,lem:comb-sym}).
  \item \emph{Thread.}  A matching matrix whose matching satisfies the
    prefix condition is factored explicitly.  In the unsymmetric case each
    matched entry is realized by its own rank-one term
    $\ell_t u_t^T$, and the prefix condition is exactly Hall's condition for
    assigning a distinct diagonal slot $t \le \min(i,j)$ to each matched pair
    $(i,j)$; a greedy assignment succeeds (\Cref{thm:lu-matching}).  In the
    symmetric case the rank-one terms $d_t \ell_t \ell_t^T$ are symmetric and
    a single off-diagonal pair cannot be realized in isolation; instead, the
    pairs are threaded into disjoint chains rooted at isolated zero indices,
    and each chain is realized by a telescoping sum of symmetric rank-one
    terms (\Cref{thm:ldlt-matching}).
\end{enumerate}
Necessity of the criteria is a short consequence of Sylvester's rank
inequality; for the symmetric case it follows directly from the unsymmetric
case since $A = LDL^T = L \, (DL^T)$ is itself an $LU$ factorization.

\paragraph{Contributions}
\begin{itemize}
  \item A unified, self-contained, constructive proof method (reduce,
    transfer, thread) that yields necessary and sufficient conditions for
    unpivoted $LU$ and unpivoted $LDL^T$ factorizations of possibly singular
    matrices over an arbitrary field, including fields of characteristic
    $2$.
  \item A necessary and sufficient condition for the existence of an
    unpivoted $LDL^T$ factorization with a truly diagonal $D$ and possibly
    singular $L$, which we believe is new, together with the corollary that
    for symmetric matrices unpivoted $LU$ and $LDL^T$ existence are
    equivalent.
  \item A self-contained account of the matching normal forms and their
    canonicity under triangular equivalence and congruence, connecting the
    existence theory to the generalized Bruhat decomposition and rank
    profile matrix~\cite{elsner_bruhat_1979,malaschonok_2010,dumas_pernet_sultan_2017}
    and to the symmetric canonical forms of Van Dooren--Strang and
    Szechtman~\cite{vandooren_strang_2014,szechtman_2007}, with short
    elimination-style proofs valid over every field.
  \item An explicit $O(n^3)$ algorithm that either produces the
    factorization or certifies its nonexistence, implemented in exact
    arithmetic and validated exhaustively against brute-force enumeration
    over small finite fields (\Cref{sec:validation}).
\end{itemize}

\paragraph{Organization}
\Cref{sec:related} surveys related work.  \Cref{sec:notation} fixes notation.
\Cref{sec:results} states the main theorems and their corollaries.
\Cref{sec:necessity} proves necessity.  \Cref{sec:normalforms} introduces
matching matrices and proves the two reduction lemmas and the canonical-form
proposition.  \Cref{sec:transfer} proves the invariance of the criteria.
\Cref{sec:combinatorics} converts the criteria into prefix conditions on the
matching.  \Cref{sec:explicit} factors the normal forms explicitly.
\Cref{sec:assembly} assembles the proofs of the main theorems and works
through examples.  \Cref{sec:validation} describes the algorithm and its
numerical validation.  The Python implementation is listed in
\Cref{sec:appendix-code}.

\section{Related work}\label{sec:related}

\paragraph{Unpivoted LU for nonsingular matrices}
The classical existence theory assumes invertibility.  If $A$ is nonsingular,
then $A = LU$ exists if and only if all leading principal submatrices are
nonsingular, in which case $L$ and $U$ are invertible and the factorization
with unit lower triangular $L$ is unique; see
Higham~\cite{higham_gaussian_2011} for a survey and
\cite{golub_matrix_2013} for textbook treatments.  Higham notes explicitly
that when a proper leading principal submatrix is singular a factorization
"may exist, but if so it is not unique," without characterizing when.  Lau
and Markham~\cite{lau_markham_1979} and Johnson, Olesky, and van den
Driessche~\cite{johnson_inherited_1989} showed that a possibly singular $A$
has a \emph{unique} unit LU factorization precisely when the $n-1$
\emph{proper} leading principal minors are nonzero; the analysis
of~\cite{johnson_inherited_1989} (inheritance of entries, fill-in) assumes
this condition throughout and does not address existence when it fails.

\paragraph{Unpivoted LU for arbitrary matrices}
The existence question for arbitrary, possibly singular $A$ was settled by
Okunev and Johnson.  Their result, circulated as a 1997 manuscript and posted
as~\cite{okunev_johnson_2005}, recently appeared
as~\cite{johnson_okunev_2025}: $A \in \F^{n \by n}$ admits $A = LU$, with no
invertibility requirement on the factors, if and only if
\begin{equation}\label{eq:oj}
  \rk(A[1:k,1:k]) + k \ \ge\ \rk(A[1:k,:]) + \rk(A[:,1:k]),
  \qquad k = 1, \ldots, n .
\end{equation}
Their sufficiency proof is an induction on "factor complements"
(generalized Schur complements), choosing the first column of $L$ and the
first row of $U$ case by case; the published version also develops
"near-LU" factorizations with almost-triangular factors when \cref{eq:oj}
fails, and re-derives $PLU$ and three-factor decompositions.
Darve~\cite{darve_lu_2026} restates \cref{eq:oj} in the nullity form used in
this paper, gives a different constructive proof based on \emph{restricted
pivoting} (only rows or columns that are zero in the active Schur complement
are moved, so that undoing the internal permutations preserves
triangularity), and develops rank-revealing structure, tight sparsity bounds
for the factors, and criteria for one-sided unit-triangular factors.  The
present paper gives a third, normal-form proof of the same criterion, whose
main advantage is that it extends verbatim to the symmetric $LDL^T$ problem.

Dopico, Johnson, and Molera~\cite{dopico_multiple_2006} study the related
question in which $L$ is required to be unit lower triangular (hence
invertible): existence is characterized through a canonical form of $A$
under left multiplication by unit lower triangular matrices (the
rank-revealing minimal canonical form), and, importantly, all LU
factorizations of a singular matrix are parametrized, quantifying their
non-uniqueness.  Our normal form differs in that the reduction acts on both
sides (or by congruence), and the middle factor --- not an upper triangular
matrix --- carries the combinatorial obstruction.
Nagarajan, Devasahayam, and Soundararajan~\cite{nagarajan_products_1999}
showed that \emph{every} square matrix over a field is a product of three
triangular matrices $LUM$; the two-factor existence question is exactly the
question of when the third factor can be dispensed with.
Jeffrey~\cite{jeffrey_lu_2010} standardizes full-rank LU factorings of
rectangular and rank-deficient matrices in the form $P_r L D^{-1} U P_c$;
since row and column permutations are built in, existence is unconditional
and the unpivoted question does not arise.  For the structured class of
M-matrices, McDonald and Schneider~\cite{mcdonald_schneider_1998} gave
graph-theoretic necessary and sufficient conditions for a singular M-matrix
to admit an unpivoted (block) LU factorization into possibly singular
M-matrix factors.  Rank-revealing LU with pivoting chosen for numerical
size, rather than algebraic existence, is developed by Miranian and
Gu~\cite{miranian_gu_2003}, and the conditioning of the triangular factors
of pivoted factorizations is analyzed by
Stewart~\cite{stewart_triangular_1997}.  General surveys of matrix
factorizations include~\cite{lu_numerical_2021}.

\paragraph{Symmetric factorizations}
For symmetric positive semidefinite matrices over $\mathbb{R}$, an unpivoted
Cholesky-type factorization always exists even in the singular case: $A =
LL^T$ with $L$ lower triangular, possibly with zero diagonal entries; see
Higham~\cite{higham_cholesky_1990} for a detailed treatment.  Our symmetric
criterion is consistent with this: positive semidefiniteness forces
$\nl(A[1:k,1:k]) = \nl(A[:,1:k])$ for all $k$, so the criterion of
\Cref{thm:ldlt} holds automatically (\Cref{cor:psd}).  For symmetric
\emph{indefinite} matrices the classical approach abandons the diagonal $D$:
with symmetric pivoting, $PAP^T = LDL^T$ always exists when $D$ is allowed
to be block diagonal with $1 \by 1$ and $2 \by 2$ blocks, as introduced by
Bunch and Parlett~\cite{bunch_parlett_1971} and made efficient by Bunch and
Kaufman~\cite{bunch_kaufman_1977}; a rank-profile-revealing variant over
arbitrary fields, still with symmetric pivoting and with $2 \by 2$
antidiagonal blocks in $D$, is given by Dumas and
Pernet~\cite{dumas_pernet_2018}.  Closest in spirit to our symmetric
reduction is the elegant paper of Van Dooren and
Strang~\cite{vandooren_strang_2014}, which factors a real symmetric matrix
as $A = L P L^T$ \emph{without pivoting}, where $P$ is a signed-matching
middle factor built from diagonal entries $0, \pm 1$ and symmetric exchange
blocks; our \Cref{lem:reduction-sym} can be read as a field-general,
weighted version of their reduction (\Cref{rem:signed}).  Neither
\cite{vandooren_strang_2014} nor \cite{dumas_pernet_2018} addresses when
the middle factor can be made \emph{truly diagonal}, which is exactly the
step from $A = MPM^T$ to $A = LDL^T$ taken in
\Cref{sec:combinatorics,sec:explicit}.  The $2 \by 2$ blocks of
Bunch--Parlett reappear in our theory in a different role: they are the
irreducible blocks of the symmetric matching normal form, i.e., the
obstruction that must be threaded through earlier zero indices rather than a
device enabled by pivoting.  We are not aware of a previously published
necessary and sufficient condition for the unpivoted diagonal-$D$ problem
with possibly singular $L$.

\paragraph{Triangular canonical forms and the Bruhat decomposition}
Our reduction lemma states that every $A$ equals $M \Pi N$ with $M, N$
invertible triangular and $\Pi$ a matching matrix, and that the matching is
an invariant of $A$.  For invertible $A$ this is a form of the Bruhat
decomposition $A = LPU$, whose permutation $P$ is uniquely determined by
$A$; see Elsner~\cite{elsner_bruhat_1979} and
Strang~\cite{strang_pnas_2010,strang_lpu_2012} for elimination-based
accounts, and~\cite{dopico_multiple_2006} for the related one-sided
canonical form.  For singular matrices the permutation becomes a partial
permutation: this is the generalized Bruhat (or LEU)
decomposition~\cite{malaschonok_2010}, and the invariant matching is the
\emph{rank profile matrix} of Dumas, Pernet, and
Sultan~\cite{dumas_pernet_sultan_2017}, with which part~(1) of our
\Cref{prop:canonical} coincides.  On the symmetric side,
Szechtman~\cite{szechtman_2007} classified symmetric matrices under
congruence by triangular groups: for $\operatorname{char} \F \neq 2$ the
canonical forms under unit-triangular congruence are weighted symmetric
sub-permutation matrices over any field, while over fields of
characteristic $2$ in which every element is a square the canonical forms
must be relaxed to include the $2 \by 2$ blocks with a nonzero corner that
appear in our \Cref{def:sym-matching}; see also Bagno and
Cherniavsky~\cite{bagno_cherniavsky_2012} for the rank-control invariants
of congruence $B$-orbits over $\mathbb{C}$.  Our
\Cref{lem:reduction-sym,prop:canonical} thus recover known canonical forms;
we include short elimination-style constructive proofs both to keep the
paper self-contained and because the constructions are exactly what the
factorization algorithm of \Cref{sec:validation} executes.

\paragraph{Positioning}
In summary: the unsymmetric criterion in \Cref{thm:lu} is due to Okunev and
Johnson~\cite{okunev_johnson_2005,johnson_okunev_2025}, with alternative
proofs and structural refinements in~\cite{darve_lu_2026}; the triangular
canonical forms we use connect to the generalized Bruhat decomposition
\cite{elsner_bruhat_1979,malaschonok_2010,dumas_pernet_sultan_2017} and, in
the symmetric case, to Van Dooren--Strang~\cite{vandooren_strang_2014} and
Szechtman~\cite{szechtman_2007}.  The contributions of the present paper
are the unified normal-form route to \emph{existence} theorems, the
symmetric $LDL^T$ criterion (\Cref{thm:ldlt}) together with the equivalence
of $LU$ and $LDL^T$ existence for symmetric matrices
(\Cref{cor:sym-lu-ldlt}), the prefix-condition combinatorics and threading
constructions of \Cref{sec:combinatorics,sec:explicit}, and the
exhaustively validated constructive algorithm.

\section{Notation and preliminaries}\label{sec:notation}

Throughout, $\F$ is an arbitrary field and all matrices have entries in
$\F$.  We use the slicing notation standard in numerical linear
algebra~\cite{golub_matrix_2013}: $A[i:j, k:l]$ is the submatrix of $A$ with
rows $i$ through $j$ and columns $k$ through $l$; a bare colon denotes the
full range, so $A[:,1:k]$ is the block of the first $k$ columns and
$A[1:k,:]$ the block of the first $k$ rows.  We write $e_i$ for the $i$th
standard basis vector, $I_n$ for the identity, and $u^T$ for the transpose.

For an $m \by p$ matrix $X$, $\rk(X)$ is its rank and
\[
  \nl(X) = p - \rk(X)
\]
is the dimension of its \emph{right} null space.  In particular, for the
three blocks entering our criteria --- each of which has exactly $k$ columns
--- we have
\begin{gather*}
  \nl(A[1:k,1:k]) = k - \rk(A[1:k,1:k]),
  \qquad
  \nl(A[:,1:k]) = k - \rk(A[:,1:k]),\\
  \nl({A[1:k,:]}^T) = k - \rk(A[1:k,:]) .
\end{gather*}

A matrix $L$ is \emph{lower triangular} if $L[i,j] = 0$ for $i < j$, and $U$
is \emph{upper triangular} if $U[i,j] = 0$ for $i > j$.  We emphasize that
triangular factors in this paper are \emph{not} assumed invertible: zeros are
allowed on their diagonals.  A triangular matrix is invertible if and only if
its diagonal entries are all nonzero, in which case its inverse is triangular
of the same kind, and products of lower (resp.\ upper) triangular matrices
are lower (resp.\ upper) triangular.  We will also use repeatedly that if
$M$ is lower triangular and $N$ is upper triangular, then
\begin{equation}\label{eq:tri-slices}
  M[1:k,:] =
  \begin{pmatrix}
    M_k & 0
  \end{pmatrix},
  \qquad
  N[:,1:k] =
  \begin{pmatrix}
    N_k \\ 0
  \end{pmatrix},
\end{equation}
where $M_k = M[1:k,1:k]$ and $N_k = N[1:k,1:k]$; if $M$ (resp.\ $N$) is
invertible then so is $M_k$ (resp.\ $N_k$), since the leading block of a
triangular matrix carries its first $k$ diagonal entries.

We record the two elementary rank facts used throughout.

\begin{lemma}[Sylvester's rank inequality]\label{lem:sylvester}
  Let $B \in \F^{m \by k}$ and $C \in \F^{k \by p}$.  Then
  \[
    \rk(B) + \rk(C) - k \ \le\ \rk(BC) \ \le\ \min\bigl(\rk(B), \rk(C)\bigr) .
  \]
\end{lemma}

\begin{proof}
  The upper bound is standard.  For the lower bound, restrict the linear map
  $x \mapsto Bx$ to $\range(C)$: its image is $\range(BC)$ and its kernel is
  contained in $\ker(B)$, so
  $\rk(C) = \dim \range(C) \le \rk(BC) + \nl(B) = \rk(BC) + k - \rk(B)$.
\end{proof}

\begin{lemma}[invariance of rank and nullity]\label{lem:rank-invariance}
  If $X$ is an $m \by p$ matrix and $G \in \F^{m \by m}$, $H \in \F^{p \by p}$
  are invertible, then $\rk(GXH) = \rk(X)$ and $\nl(GXH) = \nl(X)$.
\end{lemma}

\section{Main results}\label{sec:results}

\begin{theorem}[unpivoted LU]\label{thm:lu}
  Let $A \in \F^{n \by n}$.  The following are equivalent:
  \begin{enumerate}
    \item there exist a lower triangular $L \in \F^{n \by n}$ and an upper
      triangular $U \in \F^{n \by n}$, with zeros allowed on both diagonals,
      such that $A = LU$;
    \item for every $k = 1, \ldots, n$,
      \begin{equation}\label{eq:lu-crit}
        \nl(A[1:k,1:k]) \ \le\ \nl(A[:,1:k]) + \nl({A[1:k,:]}^T) .
      \end{equation}
  \end{enumerate}
\end{theorem}

By the rank--nullity relations of \Cref{sec:notation}, \cref{eq:lu-crit} is
equivalent to the rank form
$\rk(A[1:k,1:k]) + k \ge \rk(A[1:k,:]) + \rk(A[:,1:k])$ of Okunev and
Johnson~\cite{okunev_johnson_2005,johnson_okunev_2025}.

\begin{theorem}[unpivoted $LDL^T$]\label{thm:ldlt}
  Let $A \in \F^{n \by n}$ with $A = A^T$.  The following are equivalent:
  \begin{enumerate}
    \item there exist a lower triangular $L \in \F^{n \by n}$, with zeros
      allowed on the diagonal, and a diagonal $D \in \F^{n \by n}$ such that
      $A = LDL^T$;
    \item for every $k = 1, \ldots, n$,
      \begin{equation}\label{eq:ldlt-crit}
        \nl(A[1:k,1:k]) \ \le\ 2 \, \nl(A[:,1:k]) .
      \end{equation}
  \end{enumerate}
\end{theorem}

Equivalently, in rank form, $\rk(A[1:k,1:k]) + k \ge 2 \rk(A[:,1:k])$ for
all $k$.  Note that \cref{eq:ldlt-crit} is precisely \cref{eq:lu-crit}
specialized to a symmetric matrix: when $A = A^T$ we have
${A[1:k,:]}^T = A[:,1:k]$, so the two terms on the right-hand side of
\cref{eq:lu-crit} coincide.  This gives immediately:

\begin{corollary}[symmetric matrices: LU and $LDL^T$ are equivalent]\label{cor:sym-lu-ldlt}
  A symmetric matrix $A \in \F^{n \by n}$ admits an unpivoted $LU$
  factorization if and only if it admits an unpivoted $LDL^T$ factorization.
\end{corollary}

\begin{proof}
  If $A = LDL^T$ then $A = L \, (DL^T)$ is an LU factorization.  Conversely,
  if $A = LU$ exists then \cref{eq:lu-crit} holds by \Cref{thm:lu}, which
  for symmetric $A$ is \cref{eq:ldlt-crit}, and \Cref{thm:ldlt} produces an
  $LDL^T$ factorization.
\end{proof}

The next corollaries show that the criteria specialize to the familiar
statements in the two regimes covered by the classical theory.

\begin{corollary}[nonsingular case]\label{cor:nonsingular}
  Let $A$ be nonsingular.  Then \cref{eq:lu-crit} holds if and only if all
  leading principal submatrices $A[1:k,1:k]$ are nonsingular.  In
  particular, \Cref{thm:lu,thm:ldlt} recover the classical criteria for
  nonsingular (symmetric) matrices.
\end{corollary}

\begin{proof}
  If $A$ is nonsingular, its first $k$ columns are linearly independent and
  its first $k$ rows are linearly independent, so
  $\nl(A[:,1:k]) = \nl({A[1:k,:]}^T) = 0$.  Thus \cref{eq:lu-crit} reads
  $\nl(A[1:k,1:k]) \le 0$, i.e., $A[1:k,1:k]$ is nonsingular, for every
  $k$.  The same computation applies to \cref{eq:ldlt-crit}.
\end{proof}

\begin{corollary}[positive semidefinite case]\label{cor:psd}
  Let $\F = \mathbb{R}$ and let $A$ be symmetric positive semidefinite.
  Then \cref{eq:ldlt-crit} holds, so $A = LDL^T$ exists without pivoting;
  moreover $\nl(A[1:k,1:k]) = \nl(A[:,1:k])$ for every $k$.
\end{corollary}

\begin{proof}
  Write $A = B^T B$.  Then $A[:,1:k] = B^T \, B[:,1:k]$ and
  $A[1:k,1:k] = {B[:,1:k]}^T B[:,1:k]$.  For real $B$,
  $B[:,1:k] x = 0 \iff {B[:,1:k]}^T B[:,1:k] x = 0$ (multiply by $x^T$), so
  $\rk(A[1:k,1:k]) = \rk(B[:,1:k]) \ge \rk(A[:,1:k]) \ge
  \rk(A[1:k,1:k])$, where the last inequality holds because $A[1:k,1:k]$
  consists of rows of $A[:,1:k]$.  Hence
  $\nl(A[1:k,1:k]) = \nl(A[:,1:k]) \le 2\,\nl(A[:,1:k])$.
\end{proof}

\Cref{cor:psd} is consistent with the classical fact that a singular
positive semidefinite matrix admits an unpivoted Cholesky factorization
$A = LL^T$ with possibly zero diagonal entries in
$L$~\cite{higham_cholesky_1990}.  Indefinite symmetric matrices, by
contrast, may fail \cref{eq:ldlt-crit} --- the swap matrix $A_1$ of
\cref{eq:swap-vs-exchange} fails it at $k=1$ --- which is why the classical
unpivoted theory for indefinite matrices resorts to block-diagonal $D$ with
symmetric pivoting~\cite{bunch_parlett_1971,bunch_kaufman_1977}.

\begin{remark}[zeros on the diagonals are essential]\label{rem:zeros}
  If $L$ is required to be invertible (e.g., unit lower triangular), the
  existence conditions become strictly stronger and different in kind; see
  \cite{dopico_multiple_2006} and \cite[Sec.~3]{darve_lu_2026}.  The matrix
  $A_2$ of \cref{eq:swap-vs-exchange} satisfies
  \cref{eq:lu-crit,eq:ldlt-crit} but admits no factorization with $L$ or
  $U$ invertible.  Indeed, if $A_2 = LU$ with $L$ invertible, then
  $0 = A_2[1:2,1:2] = L[1:2,1:2] \, U[1:2,1:2]$ forces $U[1:2,1:2] = 0$,
  hence $U[:,1:2] = 0$ by \cref{eq:tri-slices} and
  $A_2[:,1:2] = L \, U[:,1:2] = 0$, contradicting $A_2[3,2] = 1$; the case
  of invertible $U$ follows by transposing the symmetric $A_2$.  Every
  factorization must instead route the exchange pair $(2,3)$ through the
  zero index $1$, forcing zero diagonal entries, as in
  \cref{eq:exchange-factorizations}.
\end{remark}

\begin{remark}[field generality]\label{rem:field}
  Both theorems hold over every field $\F$, including characteristic $2$;
  all constructions below use only the field operations, and no square
  roots, no signs, and no division by $2$.  The criteria are also
  insensitive to field extensions: ranks do not change under extension, so
  $A$ factors over $\F$ if and only if it factors over any extension of
  $\F$ (this was observed for LU in~\cite{johnson_okunev_2025}).
  A subtlety specific to characteristic $2$ is that a symmetric matrix with
  zero diagonal is alternating, and, e.g., $A_1$ over $GF(2)$ is not
  congruent to any diagonal matrix; the symmetric normal form of
  \Cref{sec:normalforms} keeps genuine $2 \by 2$ blocks for exactly this
  reason.
\end{remark}

\begin{remark}[singular $A$, singular blocks]\label{rem:singular}
  Neither theorem places any restriction on $\rk(A)$.  The nonsingular
  matrix $\bigl(\begin{smallmatrix} 0 & 1\\ 1 & 1\end{smallmatrix}\bigr)$
  fails \cref{eq:lu-crit} at $k = 1$ (its leading entry vanishes while its
  first row and column are nonzero), whereas the rank-two matrix $A_2$ of
  \cref{eq:swap-vs-exchange} satisfies both criteria.  Existence is
  governed not by how singular $A$ is, but by whether every deficiency of a
  leading block is matched by global deficiencies of the corresponding rows
  and columns.
\end{remark}

\section{Necessity}\label{sec:necessity}

Necessity of both criteria is a direct consequence of Sylvester's rank
inequality applied to the leading blocks of the factorization.

\begin{proposition}[necessity]\label{prop:necessity}
  If $A = LU$ with $L$ lower and $U$ upper triangular, then
  \cref{eq:lu-crit} holds for every $k$.  Consequently, if $A = A^T$ and
  $A = LDL^T$ with $L$ lower triangular and $D$ diagonal, then
  \cref{eq:ldlt-crit} holds for every $k$.
\end{proposition}

\begin{proof}
  Fix $k$ and partition
  \[
    L =
    \begin{pmatrix}
      L_{11} & 0\\
      L_{21} & L_{22}
    \end{pmatrix},
    \qquad
    U =
    \begin{pmatrix}
      U_{11} & U_{12}\\
      0      & U_{22}
    \end{pmatrix},
  \]
  with $L_{11}, U_{11}$ of size $k \by k$.  Reading off the leading block,
  the first $k$ columns, and the first $k$ rows of $A = LU$, and using
  \cref{eq:tri-slices},
  \[
    A[1:k,1:k] = L_{11} U_{11},
    \qquad
    A[:,1:k] =
    \begin{pmatrix}
      L_{11}\\
      L_{21}
    \end{pmatrix} U_{11},
    \qquad
    A[1:k,:] = L_{11}
    \begin{pmatrix}
      U_{11} & U_{12}
    \end{pmatrix}.
  \]
  The second identity gives $\rk(A[:,1:k]) \le \rk(U_{11})$ and the third
  gives $\rk(A[1:k,:]) \le \rk(L_{11})$.  Applying \Cref{lem:sylvester} to
  $A[1:k,1:k] = L_{11} U_{11}$,
  \[
    \rk(A[1:k,1:k])
    \ \ge\ \rk(L_{11}) + \rk(U_{11}) - k
    \ \ge\ \rk(A[1:k,:]) + \rk(A[:,1:k]) - k .
  \]
  Subtracting each side from $k$ and using $\nl(X) = k - \rk(X)$ for the
  three $k$-column blocks $A[1:k,1:k]$, $A[:,1:k]$, ${A[1:k,:]}^T$ yields
  \[
    \nl(A[1:k,1:k])
    \ \le\ 2k - \rk(A[1:k,:]) - \rk(A[:,1:k])
    \ =\ \nl({A[1:k,:]}^T) + \nl(A[:,1:k]),
  \]
  which is \cref{eq:lu-crit}.

  For the symmetric statement, observe that $A = LDL^T = L \cdot (DL^T)$ is
  an LU factorization ($DL^T$ is upper triangular), so \cref{eq:lu-crit}
  holds; since $A = A^T$ gives ${A[1:k,:]}^T = A[:,1:k]$, the right-hand
  side of \cref{eq:lu-crit} equals $2\,\nl(A[:,1:k])$, which is
  \cref{eq:ldlt-crit}.
\end{proof}

\section{Matching normal forms}\label{sec:normalforms}

The engine of both sufficiency proofs is a reduction of $A$ to a sparse
normal form by invertible triangular transformations.  The normal forms are
"matrix matchings": in the unsymmetric case a matching between row indices
and column indices, in the symmetric case a matching of the index set with
itself.

\begin{definition}[matching matrix]\label{def:matching}
  A matrix $\Pi \in \F^{n \by n}$ is a \emph{matching matrix} if every row
  and every column of $\Pi$ contains at most one nonzero entry.
  Equivalently,
  \[
    \Pi = \sum_{(i,j) \in \M} w_{ij} \, e_i e_j^T,
  \]
  where the \emph{matching} $\M = \M(\Pi) \subseteq \{1,\ldots,n\}^2$ has
  pairwise distinct first coordinates and pairwise distinct second
  coordinates, and the weights $w_{ij} \in \F$ are nonzero.  We call
  $(i,j) \in \M$ a \emph{matched pair}, a row not appearing as a first
  coordinate a \emph{zero row}, and a column not appearing as a second
  coordinate a \emph{zero column}.  When all weights equal $1$, $\Pi$ is a
  partial permutation matrix.
\end{definition}

\begin{definition}[symmetric matching matrix]\label{def:sym-matching}
  A symmetric matrix $P \in \F^{n \by n}$ is a \emph{symmetric matching
  matrix} if the index set $\{1, \ldots, n\}$ can be partitioned into
  \begin{itemize}
    \item \emph{isolated zero} indices, whose rows and columns vanish;
    \item \emph{singletons} $t$ with $P[t,t] = \gamma_t \neq 0$ and no other
      nonzero entry in row or column $t$;
    \item \emph{pairs} $(i,j)$ with $i < j$,
      \[
        P[i,j] = P[j,i] = \alpha \neq 0,
        \qquad
        P[i,i] = 0,
        \qquad
        P[j,j] = \beta \in \F ,
      \]
      and no other nonzero entries in rows and columns $i, j$.
  \end{itemize}
  Each pair block
  $\bigl(\begin{smallmatrix} 0 & \alpha\\ \alpha & \beta
  \end{smallmatrix}\bigr)$ is nonsingular (determinant $-\alpha^2 \neq 0$).
  Note the asymmetry within a pair: the \emph{earlier} index $i$ carries a
  zero diagonal entry, while the later index $j$ may carry an arbitrary
  diagonal weight $\beta$.
\end{definition}

The nonzero pattern of a symmetric matching matrix is not literally a
matching when $\beta \neq 0$ (row $j$ then carries two nonzero entries), but
its \emph{combinatorial content} is the matching that pairs $i$ with $j$ and
matches each singleton with itself; \Cref{prop:canonical} below makes this
precise.

\subsection{Reduction of an arbitrary matrix}

\begin{lemma}[triangular equivalence to a matching matrix]\label{lem:reduction-lu}
  For every $A \in \F^{n \by n}$ there exist an invertible lower triangular
  $M$, an invertible upper triangular $N$, and a matching matrix $\Pi$ such
  that
  \[
    A = M \, \Pi \, N .
  \]
  The construction below computes $M$, $\Pi$, $N$ in $O(n^3)$ field
  operations, and $M$, $N$ can be taken unit triangular.
\end{lemma}

\begin{proof}
  We induct on $k$: we construct invertible unit lower (resp.\ upper)
  triangular $M^{(k)}, N^{(k)} \in \F^{k \by k}$ and matching matrices
  $\Pi^{(k)}$ with $A[1:k,1:k] = M^{(k)} \Pi^{(k)} N^{(k)}$.  For $k = 1$
  take $M^{(1)} = N^{(1)} = (1)$ and $\Pi^{(1)} = (a_{11})$, whose matching
  is $\{(1,1)\}$ if $a_{11} \neq 0$ and empty otherwise.

  Assume the claim for $k - 1$, write $M = M^{(k-1)}$, $\Pi = \Pi^{(k-1)}$,
  $N = N^{(k-1)}$, $\M = \M(\Pi)$, and border
  \[
    A[1:k,1:k] =
    \begin{pmatrix}
      A[1:k-1,1:k-1] & a\\
      c^T & b
    \end{pmatrix},
    \qquad a, c \in \F^{k-1}, \; b \in \F .
  \]
  Setting $\widehat a := M^{-1} a$ and $\widehat c^{\,T} := c^T N^{-1}$,
  \begin{equation}\label{eq:border-lu}
    A[1:k,1:k] =
    \begin{pmatrix}
      M & 0\\
      0 & 1
    \end{pmatrix}
    \begin{pmatrix}
      \Pi & \widehat a\\
      \widehat c^{\,T} & b
    \end{pmatrix}
    \begin{pmatrix}
      N & 0\\
      0 & 1
    \end{pmatrix},
  \end{equation}
  and it remains to reduce the bordered middle matrix by further unit
  triangular factors.  The reduction has two stages.

  \emph{Stage 1: strip the matched part of the border.}
  Define $x, y \in \F^{k-1}$ supported on the matched columns and rows of
  $\Pi$:
  \[
    x_j := \widehat a_i / w_{ij},
    \qquad
    y_i := \widehat c_j / w_{ij},
    \qquad (i,j) \in \M,
  \]
  with all other components zero.  Then $\Pi x$ agrees with $\widehat a$ on
  every matched row, and $y^T \Pi$ agrees with $\widehat c^{\,T}$ on every
  matched column, so
  \begin{gather*}
    e := \widehat a - \Pi x
    \quad \text{is supported on the zero rows of } \Pi,\\
    g^T := \widehat c^{\,T} - y^T \Pi
    \quad \text{is supported on the zero columns of } \Pi .
  \end{gather*}
  Moreover $y^T e = 0$ and $g^T x = 0$, because in each product the supports
  are disjoint (matched versus zero indices).  Setting
  $f := b - y^T \Pi \, x$, a direct block multiplication gives
  \begin{equation}\label{eq:peel-lu}
    \begin{pmatrix}
      \Pi & \widehat a\\
      \widehat c^{\,T} & b
    \end{pmatrix}
    =
    \begin{pmatrix}
      I & 0\\
      y^T & 1
    \end{pmatrix}
    \begin{pmatrix}
      \Pi & e\\
      g^T & f
    \end{pmatrix}
    \begin{pmatrix}
      I & x\\
      0 & 1
    \end{pmatrix}:
  \end{equation}
  indeed the product equals
  $\bigl(\begin{smallmatrix} \Pi & \Pi x + e\\ y^T \Pi + g^T & y^T \Pi x +
  y^T e + g^T x + f \end{smallmatrix}\bigr)$, and the corner reduces to
  $y^T \Pi x + f = b$.

  \emph{Stage 2: absorb the residual border.}
  The middle matrix of \cref{eq:peel-lu} now has its border supported on
  previously zero rows and columns of $\Pi$; we convert it into at most two
  new matched pairs.

  \emph{(2a) Residual column.}
  If $e \neq 0$, let $i^\star$ be the index of its \emph{first} nonzero
  component and $\alpha := e_{i^\star}$; by Stage 1, $i^\star$ is a zero row
  of $\Pi$.  Let
  \[
    R := I + (\alpha^{-1} e - u_{i^\star}) \, u_{i^\star}^T \in
    \F^{(k-1) \by (k-1)},
  \]
  where $u_{i^\star}$ is the $i^\star$th standard basis vector.  Since the
  first nonzero of $e$ sits at $i^\star$, the vector
  $\alpha^{-1}e - u_{i^\star}$ vanishes at and above position $i^\star$, so
  $R$ is unit lower triangular.  Because row $i^\star$ of $\Pi$ is zero,
  $R \, \Pi = \Pi$; and $R \, u_{i^\star} = \alpha^{-1} e$.  Hence
  \[
    \begin{pmatrix}
      \Pi & e\\
      g^T & f
    \end{pmatrix}
    =
    \begin{pmatrix}
      R & 0\\
      0 & 1
    \end{pmatrix}
    \begin{pmatrix}
      \Pi & \alpha u_{i^\star}\\
      g^T & f
    \end{pmatrix}.
  \]
  If $e = 0$, skip this step.

  \emph{(2b) Residual row.}
  Symmetrically, if $g \neq 0$, let $j^\star$ be the first nonzero index of
  $g$ and $\gamma := g_{j^\star}$; $j^\star$ is a zero column of $\Pi$.  With
  the unit upper triangular
  $S := I + u_{j^\star} (\gamma^{-1} g - u_{j^\star})^T$ we have
  $\Pi S = \Pi$ and $u_{j^\star}^T S = \gamma^{-1} g^T$, so
  \[
    \begin{pmatrix}
      \Pi & v\\
      g^T & f
    \end{pmatrix}
    =
    \begin{pmatrix}
      \Pi & v\\
      \gamma u_{j^\star}^T & f
    \end{pmatrix}
    \begin{pmatrix}
      S & 0\\
      0 & 1
    \end{pmatrix},
    \qquad v \in \{\alpha u_{i^\star}, 0\},
  \]
  where we used that $S$ acts on the columns $1, \ldots, k-1$ only, leaving
  the border column $v$ untouched.

  \emph{(2c) Corner.}
  Let $X$ denote the current middle matrix, with border column
  $v = \alpha u_{i^\star}$ or $0$, border row $\gamma u_{j^\star}^T$ or $0$,
  and corner $f$.  Three cases:
  \begin{itemize}
    \item If $e \neq 0$: row $i^\star$ of $X$ equals $(0, \ldots, 0 \mid
      \alpha)$, so with the unit lower triangular
      $T := I_k + (f/\alpha) \, u_k u_{i^\star}^T$ we get $X = T X_0$, where
      $X_0$ is $X$ with corner $0$.
    \item Else if $g \neq 0$: column $j^\star$ of $X$ equals
      $\gamma \, u_k$, so with the unit upper triangular
      $T' := I_k + (f/\gamma) \, u_{j^\star} u_k^T$ we get $X = X_0 T'$.
    \item Else $X = \Pi \oplus (f)$.
  \end{itemize}

  Collecting the factors of Stages 1--2 into the outer factors of
  \cref{eq:border-lu} yields
  $A[1:k,1:k] = M^{(k)} \, \Pi^{(k)} \, N^{(k)}$ with unit triangular
  invertible $M^{(k)}, N^{(k)}$ and
  \[
    \M(\Pi^{(k)}) = \M(\Pi) \;\cup\;
    \begin{cases}
      \{(i^\star, k)\} & \text{if } e \neq 0 \text{ (weight } \alpha)\\
      \emptyset & \text{else}
    \end{cases}
    \;\cup\;
    \begin{cases}
      \{(k, j^\star)\} & \text{if } g \neq 0 \text{ (weight } \gamma)\\
      \{(k,k)\} & \text{if } e = 0,\ g = 0,\ f \neq 0\\
      \emptyset & \text{else.}
    \end{cases}
  \]
  This is again a matching: the new pairs occupy row $i^\star$ (previously a
  zero row), column $j^\star$ (previously a zero column), and the new index
  $k$, whose row is used by at most one of $(k, j^\star)$, $(k,k)$ and whose
  column by at most one of $(i^\star, k)$, $(k,k)$.  Taking $k = n$
  completes the induction.  Each border step costs $O(k^2)$ operations (two
  triangular solves and rank-one updates of $M$ and $N$), for $O(n^3)$ in
  total.  Only additions, multiplications, and divisions by the nonzero
  scalars $w_{ij}, \alpha, \gamma$ are used, so the construction is valid
  over every field.
\end{proof}

\begin{remark}[normalizing the weights]\label{rem:normalize-weights}
  Writing $W$ for the diagonal matrix with $W[j,j] = w_{ij}$ when column $j$
  is matched by $(i,j) \in \M$ and $W[j,j] = 1$ otherwise, we have
  $\Pi = \Pi_0 W$ with $\Pi_0$ a partial permutation matrix; absorbing $W$
  into $N$ shows that the middle factor in \Cref{lem:reduction-lu} can be
  taken to be a $0$--$1$ partial permutation matrix.  For invertible $A$
  there are no zero rows or columns, $\Pi_0$ is a permutation, and
  $A = M \Pi_0 N'$ is the Bruhat (LPU)
  decomposition~\cite{elsner_bruhat_1979,strang_lpu_2012}; for singular $A$
  it is the generalized Bruhat (LEU)
  decomposition~\cite{malaschonok_2010,dumas_pernet_sultan_2017}.
\end{remark}

\subsection{Reduction of a symmetric matrix}

Under congruence the two outer factors are transposes of one another, and
weights can no longer be normalized away (only multiplied by squares); this
is why the symmetric normal form retains the weights $\alpha$ and the
diagonal fill $\beta$.

\begin{lemma}[triangular congruence to a symmetric matching matrix]\label{lem:reduction-sym}
  For every symmetric $A \in \F^{n \by n}$ there exist an invertible lower
  triangular $M$ and a symmetric matching matrix $P$ such that
  \[
    A = M \, P \, M^T .
  \]
  The construction below computes $M$, $P$ in $O(n^3)$ field operations,
  and $M$ can be taken unit lower triangular.
\end{lemma}

\begin{proof}
  As before we induct on $k$, with base case
  $A[1:1,1:1] = (1)(a_{11})(1)$: the index $1$ is a singleton if
  $a_{11} \neq 0$ and an isolated zero otherwise.  Assume
  $A[1:k-1,1:k-1] = M P M^T$ with $M$ unit lower triangular and $P$ a
  symmetric matching matrix, and border
  \[
    A[1:k,1:k] =
    \begin{pmatrix}
      A[1:k-1,1:k-1] & a\\
      a^T & b
    \end{pmatrix},
    \qquad
    c := M^{-1} a,
  \]
  so that
  \begin{equation}\label{eq:border-sym}
    A[1:k,1:k] =
    \begin{pmatrix}
      M & 0\\
      0 & 1
    \end{pmatrix}
    \begin{pmatrix}
      P & c\\
      c^T & b
    \end{pmatrix}
    \begin{pmatrix}
      M^T & 0\\
      0 & 1
    \end{pmatrix}.
  \end{equation}

  \emph{Stage 1: strip the matched part of the border.}
  $P$ is invertible on each of its nonzero blocks, so we may choose
  $\xi \in \F^{k-1}$, supported on the singleton and pair indices, with
  $P \xi$ equal to $c$ there: for a singleton $t$ set
  $\xi_t = c_t / \gamma_t$; for a pair $(i,j)$ with weights
  $(\alpha, \beta)$ solve the $2 \by 2$ system, i.e.
  $\xi_j = c_i / \alpha$ and $\xi_i = (c_j - \beta \xi_j)/\alpha$.  Then
  \[
    e := c - P\xi
  \]
  is supported on the isolated zero indices of $P$, and $\xi^T e = 0$
  (disjoint supports).  With $f := b - \xi^T P \xi$, a direct block
  multiplication (using $\xi^T e = e^T \xi = 0$ in the corner) gives
  \begin{equation}\label{eq:peel-sym}
    \begin{pmatrix}
      P & c\\
      c^T & b
    \end{pmatrix}
    =
    \begin{pmatrix}
      I & 0\\
      \xi^T & 1
    \end{pmatrix}
    \begin{pmatrix}
      P & e\\
      e^T & f
    \end{pmatrix}
    \begin{pmatrix}
      I & \xi\\
      0 & 1
    \end{pmatrix}.
  \end{equation}

  \emph{Stage 2: absorb the residual border.}
  If $e = 0$, the middle matrix is $P \oplus (f)$: append the singleton $k$
  with weight $\gamma_k = f$ if $f \neq 0$, or the isolated zero $k$ if
  $f = 0$.  If $e \neq 0$, let $i^\star$ be its first nonzero index and
  $\alpha := e_{i^\star} \neq 0$; since $e$ is supported on isolated
  zeros, row and column $i^\star$ of $P$ vanish.  The unit lower triangular
  $R := I + (\alpha^{-1} e - u_{i^\star}) u_{i^\star}^T$ satisfies
  $R P R^T = P$ and $R u_{i^\star} = \alpha^{-1} e$, whence
  \[
    \begin{pmatrix}
      P & e\\
      e^T & f
    \end{pmatrix}
    =
    \begin{pmatrix}
      R & 0\\
      0 & 1
    \end{pmatrix}
    \begin{pmatrix}
      P & \alpha u_{i^\star}\\
      \alpha u_{i^\star}^T & f
    \end{pmatrix}
    \begin{pmatrix}
      R^T & 0\\
      0 & 1
    \end{pmatrix}.
  \]
  The middle matrix is the symmetric matching matrix obtained from $P$ by
  adding the pair $(i^\star, k)$ with weights $\alpha$ and $\beta = f$: the
  earlier index $i^\star$ (previously isolated) has zero diagonal entry, as
  required by \Cref{def:sym-matching}.  Absorbing the congruence factors of
  \cref{eq:peel-sym} and Stage~2 into the outer factors of
  \cref{eq:border-sym} completes the induction.  The cost analysis and
  field-generality are as in \Cref{lem:reduction-lu}; in particular no
  division by $2$ occurs, so the construction is valid in characteristic
  $2$.
\end{proof}

\begin{remark}[relation to known canonical forms]\label{rem:known-forms}
  Symmetric elimination without pivoting was studied by Van Dooren and
  Strang~\cite{vandooren_strang_2014}, who factor a real symmetric matrix
  as $A = LPL^T$ with $P$ built from diagonal entries $0, \pm 1$ and
  symmetric exchange blocks; \Cref{lem:reduction-sym} is a field-general,
  weighted version of that reduction.  From the canonical-form standpoint,
  Szechtman~\cite{szechtman_2007} classified symmetric matrices under
  congruence by triangular groups: for $\operatorname{char} \F \neq 2$ the
  canonical forms under unit-triangular congruence are weighted symmetric
  sub-permutation matrices, over any field, while for fields of
  characteristic $2$ (assuming every element is a square) the list of
  canonical forms necessarily includes $2 \by 2$ blocks with a nonzero
  corner --- precisely the $\beta$-entries retained in
  \Cref{def:sym-matching}.  The proof given above is included because it is
  short, uniform in the field and the characteristic, and is literally the
  algorithm implemented in \Cref{sec:validation}.
\end{remark}

\begin{remark}[signed matching matrices over $\mathbb{R}$]\label{rem:signed}
  Over $\mathbb{R}$ (more generally, over any field of characteristic
  $\neq 2$ in which every element or its negative is a square) the
  symmetric normal form can be normalized further by a diagonal-plus-shear
  congruence, applied blockwise.  For a pair block, with $i < j$,
  \[
    \begin{pmatrix}
      1 & 0\\
      -\beta/(2\alpha) & 1
    \end{pmatrix}
    \begin{pmatrix}
      0 & \alpha\\
      \alpha & \beta
    \end{pmatrix}
    \begin{pmatrix}
      1 & -\beta/(2\alpha)\\
      0 & 1
    \end{pmatrix}
    =
    \begin{pmatrix}
      0 & \alpha\\
      \alpha & 0
    \end{pmatrix},
  \]
  where the shear adds a multiple of row/column $i$ to row/column $j > i$
  and is therefore lower triangular; scaling index $j$ by $\alpha^{-1}$
  turns the block into
  $\bigl(\begin{smallmatrix} 0 & 1\\ 1 & 0 \end{smallmatrix}\bigr)$, and
  scaling a singleton $t$ by $|\gamma_t|^{-1/2}$ turns it into
  $\sgn(\gamma_t)$.  Every real symmetric matrix is thus congruent, by an
  invertible lower triangular congruence, to a \emph{signed matching
  matrix}: a symmetric matrix with entries in $\{0, \pm 1\}$, at most one
  nonzero per row and column, diagonal entries in $\{0, \pm 1\}$, and
  off-diagonal exchange blocks
  $\bigl(\begin{smallmatrix} 0 & 1\\ 1 & 0 \end{smallmatrix}\bigr)$.  The
  division by $2$ makes this normalization unavailable in characteristic
  $2$, which is why \Cref{lem:reduction-sym} retains $\beta$; over
  $GF(2)$, e.g., $\bigl(\begin{smallmatrix} 0 & 1\\ 1 & 1
  \end{smallmatrix}\bigr)$ is not congruent to
  $\bigl(\begin{smallmatrix} 0 & 1\\ 1 & 0 \end{smallmatrix}\bigr)$: the
  former is nonalternating, the latter alternating, and congruence
  preserves this distinction.
\end{remark}

\subsection{The normal form is canonical}

The matchings produced by
\Cref{lem:reduction-lu,lem:reduction-sym} are not artifacts of the
particular construction: they are invariants of $A$, computable from its
rank profile.  Part~(1) below coincides with the \emph{rank profile matrix}
of Dumas, Pernet, and Sultan~\cite{dumas_pernet_sultan_2017}; part~(2) is
the congruence analogue (cf.~\cite{bagno_cherniavsky_2012,szechtman_2007}).
For $1 \le i, j \le n$ let
\[
  r_{ij}(A) := \rk\bigl(A[1:i,1:j]\bigr),
  \qquad
  \Delta_{ij}(A) := r_{ij} - r_{i-1,j} - r_{i,j-1} + r_{i-1,j-1},
\]
with the convention $r_{0j} = r_{i0} = 0$.  We call $(i,j)$ a
\emph{rank jump} of $A$ if $\Delta_{ij}(A) = 1$; that
$\Delta_{ij} \in \{0,1\}$ always holds follows from
\Cref{lem:reduction-lu} together with part~(1) below (or
see~\cite{dumas_pernet_sultan_2017}).

\begin{proposition}[canonicity]\label{prop:canonical}
  \begin{enumerate}
    \item If $A = M \Pi N$ as in \Cref{lem:reduction-lu}, then
      \[
        r_{ij}(A) = \#\{(i',j') \in \M(\Pi) : i' \le i,\ j' \le j\}
        \quad \text{for all } i, j,
      \]
      and consequently $\M(\Pi)$ is exactly the set of rank jumps of $A$.
      In particular the matching --- and hence the zero rows and zero
      columns of $\Pi$ --- is uniquely determined by $A$ (the weights are
      not: diagonal scalings can be absorbed into $M$ or $N$).
    \item If $A = A^T = M P M^T$ as in \Cref{lem:reduction-sym}, then the
      singletons of $P$ are the diagonal rank jumps $\{t :
      \Delta_{tt}(A) = 1\}$ and the pairs of $P$ are the unordered
      off-diagonal rank jumps $\{\{i,j\} : i \neq j,\ \Delta_{ij}(A) = 1\}$;
      the partition of \Cref{def:sym-matching} is uniquely determined by
      $A$.
  \end{enumerate}
\end{proposition}

\begin{proof}
  (1) By \cref{eq:tri-slices},
  $A[1:i,1:j] = M[1:i,:] \, \Pi \, N[:,1:j] = M_i \, \Pi[1:i,1:j] \, N_j$
  with $M_i, N_j$ invertible, so $r_{ij}(A) = r_{ij}(\Pi)$ by
  \Cref{lem:rank-invariance}.  The nonzero columns of $\Pi[1:i,1:j]$ are
  the vectors $w_{i'j'} e_{i'}$ for the matched pairs with
  $i' \le i,\ j' \le j$; they occupy distinct rows, hence are linearly
  independent, giving the rank formula.  The function
  $(i,j) \mapsto \#\{(i',j') \in \M : i' \le i, j' \le j\}$ is a sum of
  corner indicators $\mathbf{1}_{\{i \ge i',\, j \ge j'\}}$, and
  $\Delta_{ij}(\mathbf{1}_{\{i \ge a,\, j \ge b\}}) = \mathbf{1}_{\{(i,j) =
  (a,b)\}}$, so the rank jumps of $A$ are exactly the matched pairs.

  (2) Congruence is the special case $N = M^T$ of a triangular
  equivalence, so as in (1), $r_{ij}(A) = r_{ij}(P)$ for all $i,j$.  The
  nonzero blocks of $P$ occupy pairwise disjoint row sets and column sets,
  so any rectangle $P[1:i,1:j]$ decomposes, after row and column
  permutations, into a direct sum of the blocks' restrictions, and its rank
  is the sum of the blocks' contributions.  A singleton $t$ contributes
  $\mathbf{1}_{\{i \ge t,\, j \ge t\}}$.  A pair $(i_0, j_0)$, $i_0 < j_0$,
  restricts to a matrix of rank $2$ when both its rows and columns are
  present, i.e.\ on the region $\{i \ge j_0,\ j \ge j_0\}$ (the two entries
  $\alpha$ are then present and the block has determinant
  $-\alpha^2 \neq 0$, regardless of $\beta$); it restricts to a matrix of
  rank $\ge 1$ exactly when at least one entry $\alpha$ is present, i.e.\ on
  the union $\{i \ge i_0,\ j \ge j_0\} \cup \{i \ge j_0,\ j \ge i_0\}$,
  whose intersection is $\{i \ge j_0,\ j \ge j_0\}$.  Writing the rank
  contribution as
  $\mathbf{1}_{\{\text{rank} \ge 1\}} + \mathbf{1}_{\{\text{rank} = 2\}}$
  and using inclusion--exclusion for the union, the contribution of the pair
  equals
  \[
    \Bigl(
    \mathbf{1}_{\{i \ge i_0,\, j \ge j_0\}}
    + \mathbf{1}_{\{i \ge j_0,\, j \ge i_0\}}
    - \mathbf{1}_{\{i \ge j_0,\, j \ge j_0\}}
    \Bigr)
    + \mathbf{1}_{\{i \ge j_0,\, j \ge j_0\}}
    =
    \mathbf{1}_{\{i \ge i_0,\, j \ge j_0\}}
    + \mathbf{1}_{\{i \ge j_0,\, j \ge i_0\}} .
  \]
  Taking mixed
  differences, the rank jumps contributed are exactly $(t,t)$ for
  singletons and $(i_0, j_0), (j_0, i_0)$ for pairs --- in particular
  $\beta$ produces no jump at $(j_0, j_0)$.  Since blocks are disjoint, the
  jump sets are disjoint, and the claim follows.
\end{proof}

\begin{remark}[compatibility of the two normal forms]\label{rem:compatibility}
  Let $A$ be symmetric.  By \Cref{prop:canonical}, the rank-jump set of $A$
  is symmetric ($\Delta_{ij} = \Delta_{ji}$), and the unsymmetric normal
  form $\Pi$ of $A$ has matching
  \[
    \M(\Pi) = \{(t,t) : t \text{ singleton of } P\}
    \;\cup\;
    \{(i,j), (j,i) : (i,j) \text{ pair of } P\} .
  \]
  Thus each $2 \by 2$ pair block of the congruence normal form corresponds
  to the \emph{two} matched pairs $(i,j)$ and $(j,i)$ of the equivalence
  normal form.  This dictionary is what makes
  \Cref{cor:sym-lu-ldlt} structurally transparent, and we will see it again
  in the combinatorial conditions of \Cref{sec:combinatorics}.
\end{remark}

\section{Invariance of the criteria}\label{sec:transfer}

\begin{lemma}[transfer]\label{lem:transfer}
  Let $M \in \F^{n \by n}$ be invertible lower triangular, $N \in
  \F^{n \by n}$ invertible upper triangular, and $B = M A N$.  Then for
  every $k = 1, \ldots, n$,
  \begin{align}
    \nl(B[1:k,1:k]) &= \nl(A[1:k,1:k]), \label{eq:transfer-lead}\\
    \nl(B[:,1:k]) &= \nl(A[:,1:k]), \label{eq:transfer-cols}\\
    \nl({B[1:k,:]}^T) &= \nl({A[1:k,:]}^T). \label{eq:transfer-rows}
  \end{align}
  Consequently $A$ satisfies \cref{eq:lu-crit} for every $k$ if and only if
  $B$ does; and, taking $N = M^T$, a symmetric $A$ satisfies
  \cref{eq:ldlt-crit} for every $k$ if and only if the congruent matrix
  $B = M A M^T$ does.
\end{lemma}

\begin{proof}
  Write $M_k = M[1:k,1:k]$ and $N_k = N[1:k,1:k]$; both are invertible
  (\Cref{sec:notation}).  Using \cref{eq:tri-slices} three times:
  \begin{itemize}
    \item $B[1:k,1:k] = M[1:k,:] \, A \, N[:,1:k]
      = M_k \, A[1:k,1:k] \, N_k$, since the zero blocks in
      \cref{eq:tri-slices} suppress the rows of $A$ beyond $k$ on the left
      and the columns beyond $k$ on the right;
    \item $B[:,1:k] = M A \, N[:,1:k] = M \, A[:,1:k] \, N_k$;
    \item $B[1:k,:] = M[1:k,:] \, A N = M_k \, A[1:k,:] \, N$, and hence
      ${B[1:k,:]}^T = N^T \, {A[1:k,:]}^T \, M_k^T$.
  \end{itemize}
  In each identity the outer factors are invertible and the two sides have
  the same number of columns, so \Cref{lem:rank-invariance} gives
  \cref{eq:transfer-lead,eq:transfer-cols,eq:transfer-rows}.  The
  congruence case follows because $M^T$ is invertible upper triangular, and
  for symmetric $A$ the quantity $\nl({A[1:k,:]}^T)$ in \cref{eq:lu-crit}
  equals $\nl(A[:,1:k])$.
\end{proof}

Combining \Cref{lem:transfer} with the reductions of
\Cref{sec:normalforms} (applied with $M^{-1}$ and $N^{-1}$, which are again
invertible triangular):

\begin{corollary}\label{cor:transfer-to-normal}
  Let $A = M \Pi N$ as in \Cref{lem:reduction-lu}.  Then $A$ satisfies
  \cref{eq:lu-crit} for every $k$ if and only if $\Pi$ does.  Let
  $A = M P M^T$ as in \Cref{lem:reduction-sym}.  Then $A$ satisfies
  \cref{eq:ldlt-crit} for every $k$ if and only if $P$ does.
\end{corollary}

\section{The criteria on the normal forms are prefix conditions}\label{sec:combinatorics}

On a matching matrix, all three nullities entering the criteria are counts
of simple combinatorial objects, and the criteria collapse to conditions on
prefixes of the index set.

\subsection{The unsymmetric normal form}

Let $\Pi$ be a matching matrix with matching $\M$.  Relative to the cut
after index $k$, classify the matched pairs and define
\begin{align*}
  \iota_k &:= \#\{(i,j) \in \M : i \le k \text{ and } j \le k\}
  && \text{(pairs inside the leading block)},\\
  \chi^R_k &:= \#\{(i,j) \in \M : i \le k < j\}
  && \text{(pairs crossing the cut row-first)},\\
  \chi^C_k &:= \#\{(i,j) \in \M : j \le k < i\}
  && \text{(pairs crossing the cut column-first)},\\
  r^0_k &:= \#\{\text{zero rows} \le k\},
  \qquad
  c^0_k := \#\{\text{zero columns} \le k\}. \hspace{-12em}
\end{align*}

\begin{lemma}[nullities of a matching matrix]\label{lem:null-matching}
  For every $k$,
  \begin{equation}\label{eq:null-matching}
    \nl(\Pi[1:k,1:k]) = r^0_k + \chi^R_k = c^0_k + \chi^C_k,
    \qquad
    \nl(\Pi[:,1:k]) = c^0_k,
    \qquad
    \nl({\Pi[1:k,:]}^T) = r^0_k .
  \end{equation}
  In particular the \emph{bookkeeping identity}
  $r^0_k + \chi^R_k = c^0_k + \chi^C_k$ holds.
\end{lemma}

\begin{proof}
  The nonzero columns of the $n \by k$ block $\Pi[:,1:k]$ are the vectors
  $w_{ij} e_i$ over the matched pairs with $j \le k$; they occupy distinct
  rows and are therefore linearly independent, while the remaining $c^0_k$
  columns vanish.  Hence $\rk(\Pi[:,1:k]) = k - c^0_k$, giving the second
  formula, and the third follows by applying the same argument to $\Pi^T$,
  whose matching is the transpose of $\M$.

  The leading block $\Pi[1:k,1:k]$ is itself a matching matrix whose
  matched pairs are the $\iota_k$ pairs inside the block, so
  $\rk(\Pi[1:k,1:k]) = \iota_k$ (as in \Cref{prop:canonical}, its nonzero
  columns are independent).  Classify the $k$ row indices $\le k$: each is
  a zero row, or belongs to a pair with $j \le k$, or to a pair with
  $j > k$; the classes are disjoint and exhaust, so
  $k = r^0_k + \iota_k + \chi^R_k$ and
  $\nl(\Pi[1:k,1:k]) = k - \iota_k = r^0_k + \chi^R_k$.  Classifying the
  column indices instead gives $k = c^0_k + \iota_k + \chi^C_k$ and the
  alternative expression, whence the bookkeeping identity.
\end{proof}

\begin{lemma}[prefix form of the LU criterion]\label{lem:comb-lu}
  For a matching matrix $\Pi$, criterion \cref{eq:lu-crit} at index $k$ is
  equivalent to each of
  \[
    \text{(a) } \chi^R_k \le c^0_k,
    \qquad
    \text{(b) } \chi^C_k \le r^0_k,
    \qquad
    \text{(c) } \nu_k := \#\{(i,j) \in \M : \min(i,j) \le k\} \ \le\ k .
  \]
\end{lemma}

\begin{proof}
  Substituting \cref{eq:null-matching} into \cref{eq:lu-crit}: using the
  first expression for the leading nullity,
  $r^0_k + \chi^R_k \le c^0_k + r^0_k$, i.e.\ (a); using the second,
  $c^0_k + \chi^C_k \le c^0_k + r^0_k$, i.e.\ (b).  By the bookkeeping
  identity, (a) and (b) are equivalent (they have equal slack).  For (c):
  a pair has $\min(i,j) \le k$ iff it is counted by $\iota_k$, $\chi^R_k$,
  or $\chi^C_k$, so $\nu_k = \iota_k + \chi^R_k + \chi^C_k$; substituting
  $\iota_k = k - r^0_k - \chi^R_k$ from the proof of
  \Cref{lem:null-matching} gives
  $\nu_k \le k \iff \chi^C_k \le r^0_k$, which is (b).
\end{proof}

\subsection{The symmetric normal form}

Let $P$ be a symmetric matching matrix.  Define
\[
  \zeta_k := \#\{\text{isolated zero indices} \le k\},
  \qquad
  \chi_k := \#\{\text{pairs } (i,j) \text{ of } P : i \le k < j\} .
\]

\begin{lemma}[nullities of a symmetric matching matrix]\label{lem:null-sym-matching}
  For every $k$,
  \[
    \nl(P[1:k,1:k]) = \zeta_k + \chi_k,
    \qquad
    \nl(P[:,1:k]) = \zeta_k .
  \]
\end{lemma}

\begin{proof}
  The columns $t \le k$ of the $n \by k$ block $P[:,1:k]$ are: zero for
  isolated $t$; $\gamma_t e_t$ for singletons; for a pair $(i,j)$ of $P$,
  column $i$ of $P$ equals $\alpha e_j$ and column $j$ equals
  $\alpha e_i + \beta e_j$ (whichever of the two has index $\le k$ is
  included).  Distinct blocks occupy disjoint row sets, and within a pair
  with both $i, j \le k$ the two columns $\alpha e_j$ and
  $\alpha e_i + \beta e_j$ are independent (restricted to rows $\{i,j\}$
  they form the pair block, of determinant $-\alpha^2$).  Hence the nonzero
  columns of $P[:,1:k]$ are independent, exactly the isolated columns
  vanish, and $\nl(P[:,1:k]) = \zeta_k$.

  The leading block $P[1:k,1:k]$ is block diagonal (up to symmetric
  permutation) with blocks: zeros at isolated indices ($\zeta_k$ of them);
  $(\gamma_t)$ at singletons (rank $1$); full pair blocks when
  $j \le k$ (rank $2$); and, for each crossing pair $i \le k < j$, the
  single row/column $i$ whose entries within the block all vanish
  (its only nonzero, $\alpha$, sits in column $j > k$) --- contributing one
  null direction each.  Hence
  $\nl(P[1:k,1:k]) = \zeta_k + \chi_k$.
\end{proof}

\begin{lemma}[prefix form of the $LDL^T$ criterion]\label{lem:comb-sym}
  For a symmetric matching matrix $P$, criterion \cref{eq:ldlt-crit} at
  index $k$ is equivalent to
  \begin{equation}\label{eq:prefix-sym}
    \chi_k \ \le\ \zeta_k .
  \end{equation}
\end{lemma}

\begin{proof}
  Substitute \Cref{lem:null-sym-matching} into \cref{eq:ldlt-crit}:
  $\zeta_k + \chi_k \le 2 \zeta_k$.
\end{proof}

\begin{remark}\label{rem:comb-dictionary}
  Under the dictionary of \Cref{rem:compatibility} --- a pair $(i,j)$ of
  $P$ corresponds to the two matched pairs $(i,j)$ and $(j,i)$ of the
  unsymmetric normal form $\Pi$ of the same symmetric $A$, an isolated zero
  to a common zero row and column --- we have
  $\chi^R_k = \chi^C_k = \chi_k$ and $r^0_k = c^0_k = \zeta_k$, and
  conditions (a), (b) of \Cref{lem:comb-lu} both become
  \cref{eq:prefix-sym}.  The combinatorial obstruction is thus literally
  the same in the two problems.
\end{remark}

\section{Explicit factorization of the normal forms}\label{sec:explicit}

We now factor the normal forms under the prefix conditions of
\Cref{sec:combinatorics}.  Both constructions build the factorization as a
sum of rank-one terms indexed by \emph{slots} $t \in \{1, \ldots, n\}$: for
$A = LU$,
\begin{equation}\label{eq:rank-one-lu}
  LU = \sum_{t=1}^{n} \ell_t \, u_t^T,
\end{equation}
where $\ell_t$ is column $t$ of $L$ (supported on rows $\ge t$) and $u_t^T$
is row $t$ of $U$ (supported on columns $\ge t$); for $A = LDL^T$,
\begin{equation}\label{eq:rank-one-ldlt}
  LDL^T = \sum_{t=1}^{n} d_t \, \ell_t \ell_t^T .
\end{equation}
The difference between the two problems is now plain: in
\cref{eq:rank-one-lu} each term is an arbitrary rank-one matrix with the
required support, whereas in \cref{eq:rank-one-ldlt} each term is
\emph{symmetric} rank-one.  A single off-diagonal exchange block
$\alpha(e_i e_j^T + e_j e_i^T)$ has rank $2$ and zero diagonal, and cannot
equal a sum of terms each supported in rows and columns $\ge i$ unless
partial sums are allowed to cancel (over any field: the $(i,i)$ entry
forces $d_t \ell_t[i]^2 = 0$ for every term while the $(i,j)$ entry needs
some $d_t \ell_t[i] \neq 0$); this forces the telescoping chains of
\Cref{thm:ldlt-matching}.

\subsection{LU of a matching matrix: greedy slot assignment}

\begin{theorem}\label{thm:lu-matching}
  Let $\Pi$ be a matching matrix satisfying the prefix condition
  $\nu_k \le k$ for all $k$ (\Cref{lem:comb-lu}(c)).  Then there is an
  injection
  \begin{equation}\label{eq:slot-injection}
    \sigma : \M(\Pi) \hookrightarrow \{1, \ldots, n\},
    \qquad
    \sigma(i,j) \le \min(i,j) \quad \text{for every } (i,j) \in \M(\Pi),
  \end{equation}
  computable by a greedy pass, and setting, for each $(i,j) \in \M$ with
  $t = \sigma(i,j)$,
  \[
    \ell_t := e_i,
    \qquad
    u_t := w_{ij} \, e_j,
  \]
  and $\ell_t = u_t = 0$ for slots $t \notin \range(\sigma)$, the matrices
  $L = (\ell_1, \ldots, \ell_n)$ and $U$ with rows $u_1^T, \ldots, u_n^T$
  are lower and upper triangular, respectively, and $L U = \Pi$.
\end{theorem}

\begin{proof}
  \emph{The assignment suffices.}
  Given $\sigma$ as in \cref{eq:slot-injection}, triangularity is immediate:
  $\ell_t = e_i$ has its support at row $i \ge \min(i,j) \ge t$, and $u_t$
  at column $j \ge t$.  Since $\sigma$ is injective, each matched pair
  contributes its own term in \cref{eq:rank-one-lu} and empty slots
  contribute zero, so
  $L U = \sum_{(i,j) \in \M} e_i \, (w_{ij} e_j)^T = \Pi$.  Note that both
  diagonals may carry zeros: slot $t$ has $L[t,t] = 1$ only if
  $i = t$, and $U[t,t] \neq 0$ only if $j = t$.

  \emph{The assignment exists.}
  Process the pairs in order of nondecreasing $m(i,j) := \min(i,j)$, ties
  broken arbitrarily, maintaining the set of free slots (initially
  $\{1, \ldots, n\}$); assign to each pair the largest free slot
  $t \le m(i,j)$ (the argument below is valid for \emph{any} choice of
  free slot $\le m(i,j)$; we fix the largest for definiteness).  Suppose
  the greedy fails at a pair with bound
  $m := m(i,j)$: all slots $1, \ldots, m$ are occupied.  Every occupied
  slot $s \le m$ was assigned to a previously processed pair, whose bound
  was $\le m$ (bounds are nondecreasing along the pass).  Together with the
  current pair, this exhibits $m + 1$ pairs with $\min \le m$, i.e.\
  $\nu_m \ge m + 1$, contradicting the prefix condition.  Hence the greedy
  assigns every pair and produces $\sigma$.
\end{proof}

\begin{remark}[which pairs need rerouting]\label{rem:rerouting}
  Order the pairs by their bound $m(i,j)$.  A diagonal pair $(t,t)$ can
  take its own slot $t$; a pair $(i,j)$ with $i < j$ (an "above" pair)
  would naturally take slot $i$, and a "below" pair $(i,j)$ with $i > j$
  slot $j$.  Conflicts occur exactly at \emph{crossing indices} $t$ that
  are simultaneously the row of an above pair $(t, j)$ and the column of a
  below pair $(i, t)$: one of the two must be rerouted to an earlier free
  slot, whose availability is exactly what the prefix condition
  guarantees --- in the leanest case the free slot comes from an index that
  is both a zero row and a zero column.  The exchange matrix $A_2$ of
  \cref{eq:swap-vs-exchange} is the minimal example: index $2$ is crossing,
  and the pair $(2,3)$ (or $(3,2)$) is rerouted through slot $1$, whose row
  and column are zero.  This is the precise sense in which zero diagonal
  entries are a \emph{resource} for unpivoted factorization.
\end{remark}

\subsection{$LDL^T$ of a symmetric matching matrix: threads}

Throughout this subsection $P$ is a symmetric matching matrix satisfying
$\chi_k \le \zeta_k$ for all $k$ (\Cref{lem:comb-sym}).

\begin{lemma}[thread partition]\label{lem:threads}
  The pairs of $P$ can be partitioned into \emph{threads}
  \begin{equation}\label{eq:thread}
    r < i_1 < j_1 < i_2 < j_2 < \cdots < i_m < j_m,
  \end{equation}
  where $r$ is an isolated zero index, $(i_q, j_q)$ are pairs of $P$, and
  distinct threads use disjoint sets of indices.
\end{lemma}

\begin{proof}
  Scan the indices from left to right, maintaining a set of \emph{free
  threads}.  When an isolated zero index $r$ is encountered, create a new
  free thread rooted at $r$.  When a pair $(i,j)$ starts at index $i$,
  attach it to any free thread, which becomes busy until index $j$; when
  the scan reaches $j$, the thread becomes free again.  (Each index belongs
  to at most one block, so no two events collide at the same index.)  It
  remains to show a free thread is always available when a pair starts at
  $i$.  The threads created before index $i$ number $\zeta_{i-1} = \zeta_i$
  ($i$ is a pair index, not an isolated zero).  The threads busy at that
  moment are those attached to pairs $(i', j')$ with $i' < i < j'$;
  together with the new pair, these are exactly the pairs counted by
  $\chi_i = \#\{(i', j') : i' \le i < j'\}$, so at most $\chi_i - 1$
  threads are busy when the new pair arrives.  Since $\chi_i \le \zeta_i$,
  at least $\zeta_i - (\chi_i - 1) \ge 1$ of the created threads are
  free.  By construction
  each thread's blocks are attached after its root and after the previous
  block's end, giving the interleaving \cref{eq:thread}, and every index is
  used by at most one thread.
\end{proof}

\begin{lemma}[one thread telescopes]\label{lem:telescope}
  Fix a thread \cref{eq:thread} and write $\alpha_q \neq 0$ and $\beta_q$
  for the weights of the pair $(i_q, j_q)$.  Define, backward for
  $q = m, m-1, \ldots, 1$:
  \[
    \delta_{m+1} := 1,
    \quad
    s_{m+1} := 0 \in \F^n,
    \qquad
    v_q :=
    \begin{cases}
      0 & \beta_q \neq 0\\
      1 & \beta_q = 0
    \end{cases},
    \quad
    \delta_q := \delta_{q+1} v_q^2 - \beta_q,
    \quad
    s_q := -\frac{\alpha_q}{\delta_q} e_{i_q}
    + \frac{\delta_{q+1} v_q}{\delta_q} s_{q+1} .
  \]
  Each $\delta_q$ is nonzero ($\delta_q = -\beta_q \neq 0$ in the first
  case, $\delta_q = \delta_{q+1}$ in the second).  Define the columns and
  diagonal entries
  \[
    \ell_r := s_1, \quad d_r := \delta_1;
    \qquad
    \ell_{i_q} := s_q + e_{j_q}, \quad d_{i_q} := -\delta_q;
    \qquad
    \ell_{j_q} := s_{q+1} + v_q e_{j_q}, \quad d_{j_q} := \delta_{q+1} .
  \]
  Then each $\ell_t$ is supported on rows $\ge t$, and
  \begin{equation}\label{eq:thread-sum}
    d_r \ell_r \ell_r^T
    + \sum_{q=1}^{m}
    \bigl( d_{i_q} \ell_{i_q} \ell_{i_q}^T + d_{j_q} \ell_{j_q}
    \ell_{j_q}^T \bigr)
    =
    \sum_{q=1}^{m}
    \Bigl( \alpha_q \bigl( e_{i_q} e_{j_q}^T + e_{j_q} e_{i_q}^T \bigr)
    + \beta_q \, e_{j_q} e_{j_q}^T \Bigr),
  \end{equation}
  i.e.\ the thread's columns reproduce exactly the thread's blocks of $P$.
\end{lemma}

\begin{proof}
  \emph{Support.}  By downward induction, $s_q$ is supported on
  $\{i_q, i_{q+1}, \ldots, i_m\}$.  Hence $\ell_r = s_1$ is supported on
  indices $\ge i_1 > r$; $\ell_{i_q}$ on
  $\{i_q, \ldots, i_m\} \cup \{j_q\}$, all $\ge i_q$; and $\ell_{j_q}$ on
  $\{i_{q+1}, \ldots, i_m\} \cup \{j_q\}$, all $\ge j_q$ since
  $i_{q+1} > j_q$.

  \emph{Telescoping.}  The recursion gives the two identities
  \begin{equation}\label{eq:thread-identities}
    -\delta_q s_q + \delta_{q+1} v_q \, s_{q+1} = \alpha_q e_{i_q},
    \qquad
    \delta_{q+1} v_q^2 - \delta_q = \beta_q .
  \end{equation}
  Expanding the symmetric rank-one terms of one pair,
  \begin{align*}
    d_{i_q} \ell_{i_q} \ell_{i_q}^T + d_{j_q} \ell_{j_q} \ell_{j_q}^T
    &= -\delta_q (s_q + e_{j_q})(s_q + e_{j_q})^T
    + \delta_{q+1} (s_{q+1} + v_q e_{j_q})(s_{q+1} + v_q e_{j_q})^T \\
    &= \bigl( -\delta_q s_q s_q^T + \delta_{q+1} s_{q+1} s_{q+1}^T \bigr)
    + \bigl( -\delta_q s_q + \delta_{q+1} v_q s_{q+1} \bigr) e_{j_q}^T \\
    &\qquad
    + e_{j_q} \bigl( -\delta_q s_q + \delta_{q+1} v_q s_{q+1} \bigr)^T
    + \bigl( \delta_{q+1} v_q^2 - \delta_q \bigr) e_{j_q} e_{j_q}^T \\
    &= \bigl( -\delta_q s_q s_q^T + \delta_{q+1} s_{q+1} s_{q+1}^T \bigr)
    + \alpha_q \bigl( e_{i_q} e_{j_q}^T + e_{j_q} e_{i_q}^T \bigr)
    + \beta_q \, e_{j_q} e_{j_q}^T,
  \end{align*}
  using \cref{eq:thread-identities}.  Summing over $q = 1, \ldots, m$, the
  first brackets telescope to
  $-\delta_1 s_1 s_1^T + \delta_{m+1} s_{m+1} s_{m+1}^T
  = -\delta_1 s_1 s_1^T$, which is cancelled by the root term
  $d_r \ell_r \ell_r^T = \delta_1 s_1 s_1^T$.  This proves
  \cref{eq:thread-sum}.
\end{proof}

\begin{theorem}\label{thm:ldlt-matching}
  Let $P$ be a symmetric matching matrix with $\chi_k \le \zeta_k$ for all
  $k$.  Then there are a lower triangular $S$ and a diagonal $E$ with
  \[
    P = S E S^T .
  \]
\end{theorem}

\begin{proof}
  Partition the pairs into threads (\Cref{lem:threads}) and apply
  \Cref{lem:telescope} to each thread, obtaining columns
  $\ell_t$ and weights $d_t$ for every index $t$ in a thread.  For each
  singleton $t$ set $\ell_t := e_t$ and $d_t := \gamma_t$; for isolated
  zero indices not rooting a thread set $\ell_t := 0$, $d_t := 0$.  Let
  $S := (\ell_1, \ldots, \ell_n)$ and $E := \operatorname{diag}(d_1,
  \ldots, d_n)$.  Every $\ell_t$ is supported on rows $\ge t$, so $S$ is
  lower triangular.  Distinct threads and singletons use disjoint index
  sets, so summing \cref{eq:thread-sum} over threads and adding the
  singleton terms $\gamma_t e_t e_t^T$ reproduces every block of $P$ and
  nothing else:
  $S E S^T = \sum_t d_t \ell_t \ell_t^T = P$.
\end{proof}

\begin{remark}[the real signed case]\label{rem:signed-threads}
  Over $\mathbb{R}$, after the normalization of \Cref{rem:signed}
  ($\alpha_q = 1$, $\beta_q = 0$), the recursion of \Cref{lem:telescope}
  simplifies to $v_q = 1$, $\delta_q = 1$, and
  $s_q = s_{q+1} - e_{i_q} = -(e_{i_q} + e_{i_{q+1}} + \cdots + e_{i_m})$,
  so the thread columns are signed partial sums and
  \[
    d_r = 1, \qquad d_{i_q} = -1, \qquad d_{j_q} = 1 :
  \]
  the diagonal alternates along each thread, and $D$ can be taken with
  entries in $\{0, \pm 1\}$.  The factorization
  \cref{eq:exchange-factorizations} of $A_2$ is exactly the one-thread case
  $r = 1 < i_1 = 2 < j_1 = 3$, obtained by seeding the recursion with
  $\delta_{m+1} = -1$ instead of $1$ (valid whenever all $\beta_q = 0$),
  which flips the sign of every $\delta_q$ and $s_q$.
\end{remark}

\section{Proofs of the main theorems, and worked examples}\label{sec:assembly}

\begin{proof}[Proof of \Cref{thm:lu}]
  Necessity is \Cref{prop:necessity}.  For sufficiency, assume $A$
  satisfies \cref{eq:lu-crit} for all $k$.  By \Cref{lem:reduction-lu},
  $A = M \Pi N$ with $M, N$ invertible triangular and $\Pi$ a matching
  matrix.  By \Cref{cor:transfer-to-normal}, $\Pi$ satisfies
  \cref{eq:lu-crit}, hence by \Cref{lem:comb-lu} the prefix condition
  $\nu_k \le k$ for all $k$.  By \Cref{thm:lu-matching} there are lower
  triangular $L_\Pi$ and upper triangular $U_\Pi$ with
  $L_\Pi U_\Pi = \Pi$.  Then
  \[
    A = M \, (L_\Pi U_\Pi) \, N = (M L_\Pi)(U_\Pi N),
  \]
  and $M L_\Pi$ is lower triangular, $U_\Pi N$ upper triangular.
\end{proof}

\begin{proof}[Proof of \Cref{thm:ldlt}]
  Necessity is \Cref{prop:necessity}.  For sufficiency, assume the
  symmetric $A$ satisfies \cref{eq:ldlt-crit} for all $k$.  By
  \Cref{lem:reduction-sym}, $A = M P M^T$ with $M$ invertible lower
  triangular and $P$ a symmetric matching matrix.  By
  \Cref{cor:transfer-to-normal}, $P$ satisfies \cref{eq:ldlt-crit}, hence
  by \Cref{lem:comb-sym} the prefix condition $\chi_k \le \zeta_k$ for all
  $k$.  By \Cref{thm:ldlt-matching}, $P = S E S^T$ with $S$ lower
  triangular and $E$ diagonal.  Then
  \[
    A = M S \, E \, S^T M^T = (MS) \, E \, (MS)^T,
  \]
  and $MS$ is lower triangular.
\end{proof}

\begin{example}[the exchange matrix, both ways]\label{ex:exchange}
  The matrix $A_2$ of \cref{eq:swap-vs-exchange} is already a normal form.
  As a matching matrix, $\M = \{(2,3), (3,2)\}$ with unit weights; row $1$
  and column $1$ are zero.  The prefix counts are $\nu_1 = 0 \le 1$ and
  $\nu_2 = 2 \le 2$ (both pairs have $\min = 2$), so \Cref{thm:lu-matching}
  applies: the greedy assigns slot $2$ to $(2,3)$ and reroutes $(3,2)$ to
  slot $1$, giving $\ell_2 = e_2$, $u_2 = e_3$, $\ell_1 = e_3$,
  $u_1 = e_2$, which is precisely the LU factorization displayed in
  \cref{eq:exchange-factorizations}.  Index $2$ is a crossing index in the
  sense of \Cref{rem:rerouting}, and slot $1$ --- a zero row \emph{and}
  zero column --- is the resource that resolves it.

  As a symmetric matching matrix, index $1$ is an isolated zero and
  $(2,3)$ is a pair with $\alpha = 1$, $\beta = 0$; the prefix condition
  reads $\chi_2 = 1 \le \zeta_2 = 1$.  There is a single thread
  $1 < 2 < 3$, and the recursion of \Cref{lem:telescope} gives
  ($v_1 = 1$, $\delta_1 = \delta_2 = 1$, $s_1 = -e_2$)
  \[
    L =
    \begin{pmatrix}
      0 & 0 & 0\\
      -1 & -1 & 0\\
      0 & 1 & 1
    \end{pmatrix},
    \qquad
    D = \operatorname{diag}(1, -1, 1),
    \qquad
    L D L^T = A_2 ,
  \]
  the sign-flipped variant of \cref{eq:exchange-factorizations}.  In
  characteristic $2$ the same formulas read $L D L^T$ with
  $D = \operatorname{diag}(1,1,1)$ and $-1 = 1$, and remain valid: no step
  of the construction used signs.
\end{example}

\begin{example}[a genuine crossing]\label{ex:crossing}
  Let $n = 4$ and
  \[
    \Pi =
    \begin{pmatrix}
      0 & 0 & 0 & 0\\
      0 & 0 & 0 & 1\\
      0 & 0 & 0 & 0\\
      0 & 1 & 0 & 0
    \end{pmatrix},
    \qquad
    \M = \{(2,4), (4,2)\},
  \]
  with rows and columns $1, 3$ zero.  Both pairs have $\min = 2$, so
  $\nu_2 = 2 \le 2$; condition (a) of \Cref{lem:comb-lu} reads
  $\chi^R_2 = 1 \le c^0_2 = 1$ and $\chi^R_3 = 1 \le c^0_3 = 2$.  The
  greedy gives slot $2$ to $(2,4)$ and slot $1$ to $(4,2)$:
  \[
    L =
    \begin{pmatrix}
      0 & 0 & 0 & 0\\
      0 & 1 & 0 & 0\\
      0 & 0 & 0 & 0\\
      1 & 0 & 0 & 0
    \end{pmatrix},
    \qquad
    U =
    \begin{pmatrix}
      0 & 1 & 0 & 0\\
      0 & 0 & 0 & 1\\
      0 & 0 & 0 & 0\\
      0 & 0 & 0 & 0
    \end{pmatrix},
    \qquad
    LU = \Pi .
  \]
  Viewed symmetrically, $\Pi$ is the symmetric matching matrix with
  isolated zeros $\{1,3\}$ and the single pair $(2,4)$; the thread
  $1 < 2 < 4$ produces the corresponding $LDL^T$.
\end{example}

\begin{example}[ordering, not rank, is the obstruction]\label{ex:char2}
  Over $GF(2)$, the nonsingular symmetric matrix
  $A = \bigl(\begin{smallmatrix} 0 & 1\\ 1 & 1 \end{smallmatrix}\bigr)$
  fails both criteria at $k = 1$: $\nl(A[1:1,1:1]) = 1$ while the first
  column and row of $A$ are nonzero.  So $A$ has no unpivoted $LU$ or
  $LDL^T$ factorization.  Yet $A$ \emph{is} congruent to the identity by a
  full (non-triangular) congruence:
  with
  $T = \bigl(\begin{smallmatrix} 1 & 1\\ 0 & 1 \end{smallmatrix}\bigr)$,
  $T A T^T = I_2$.  Unpivoted existence is thus a statement about the
  \emph{ordering} of the index set, not about rank or congruence class ---
  consistent with the motivation of preserving a causal or temporal
  ordering discussed in \Cref{sec:intro}.
\end{example}

\section{Algorithm and numerical validation}\label{sec:validation}

\subsection{The algorithm}

The proofs of \Cref{thm:lu,thm:ldlt} are constructive and assemble into
the following $O(n^3)$ algorithm, implemented in
\Cref{sec:appendix-code}.  Given $A$:
\begin{enumerate}
  \item \emph{Reduce} (cost $O(n^3)$): run the bordered induction of
    \Cref{lem:reduction-lu} (resp.\ \Cref{lem:reduction-sym}) to obtain
    $A = M \Pi N$ (resp.\ $A = M P M^T$).  Each border step performs two
    triangular solves against the current $M$ and $N$ and a constant
    number of rank-one updates.
  \item \emph{Thread} (cost $O(n \log n)$): run the greedy slot assignment
    of \Cref{thm:lu-matching} (resp.\ the thread scan of
    \Cref{lem:threads} and the backward recursion of
    \Cref{lem:telescope}) on the normal form.  If the greedy runs out of
    slots (resp.\ threads), report failure: by
    \Cref{lem:comb-lu,lem:comb-sym,cor:transfer-to-normal} and
    \Cref{prop:necessity}, the index $k$ at which this happens certifies
    that no factorization exists.
  \item \emph{Assemble} (cost $O(n^3)$): return
    $L = M L_\Pi$, $U = U_\Pi N$ (resp.\ $L = M S$, $D = E$).
\end{enumerate}
The algorithm uses only field operations, and division only by provably
nonzero quantities, so it runs unchanged over $\mathbb{Q}$, over $GF(p)$,
or over any computable field.  Like the constructions
of~\cite{okunev_johnson_2005,darve_lu_2026}, it makes no attempt to select
pivots by magnitude, and is therefore not backward stable over
$\mathbb{R}$ in floating-point arithmetic; stabilizing the construction is
an open problem (\Cref{sec:conclusion}).

\subsection{Validation methodology and results}

We validate every theorem, lemma, and construction of this paper with
exact arithmetic (rational numbers via Python's \texttt{fractions}, and
prime fields $GF(2)$, $GF(3)$), so no floating-point tolerance can mask an
error.  The driver \texttt{validate.py} (see \Cref{sec:appendix-code})
runs four independent layers of checks.
\begin{enumerate}
  \item \emph{Criterion $\Leftrightarrow$ construction, with exact
    factors.}  For each test matrix, the nullity criterion
    (\cref{eq:lu-crit} or \cref{eq:ldlt-crit}) is evaluated directly from
    exact ranks, and the constructive algorithm is run.  We assert that
    the construction succeeds exactly when the criterion holds and that,
    on success, the factors are triangular (diagonal for $D$) and multiply
    back to $A$ \emph{exactly}.
  \item \emph{Brute force as independent ground truth.}  Over $GF(2)$ and
    $GF(3)$ for small $n$, we enumerate \emph{all} products $LU$ (resp.\
    $LDL^T$) over all triangular $L$, $U$ (resp.\ all $L$ and diagonal
    $D$), obtaining the exact set of factorizable matrices with no
    reference to any result of this paper, and compare it with the
    criterion on every matrix (resp.\ every symmetric matrix) of that
    size.  This independently confirms both necessity and sufficiency.
  \item \emph{Canonicity.}  For random rational matrices, the matching
    computed by the reduction is compared with the rank-jump set of
    \Cref{prop:canonical}, computed from $n^2$ independent exact rank
    evaluations.
  \item \emph{The symmetric corollary.}  For every symmetric test matrix,
    we assert that \cref{eq:lu-crit} and \cref{eq:ldlt-crit} agree, as
    \Cref{cor:sym-lu-ldlt} requires.
\end{enumerate}
Random rational matrices are drawn from four structural families (dense;
sparse; exact low rank $BC$ with $\rk \le n$; and triangularly disguised
sparse cores $T_1 S T_2$) to exercise deficient leading blocks, deficient
row/column blocks, and crossing patterns.

\Cref{tab:validation} reports the results: $31{,}639$ cases, zero
failures, about $7$ seconds in pure Python.  The exhaustive rows are
complete censuses of the corresponding matrix classes: for example, every
one of the $512$ matrices in $GF(2)^{3 \by 3}$ is compared against the
brute-force set of all $64 \times 64$ products of triangular matrices, and
every one of the $1024$ symmetric matrices in $GF(2)^{4 \by 4}$ against
all $2^{10} \times 2^4$ products $L D L^T$.

\begin{table}[htbp]
  \centering
  \begin{tabular}{lrr}
    \toprule
    Test & Cases & Failures \\
    \midrule
    Special matrices ($\mathbb{Q}$ and $GF(2)$) & 18 & 0 \\
    Random LU over $\mathbb{Q}$, $n \le 7$ (1435 satisfying \cref{eq:lu-crit}) & 2000 & 0 \\
    Random $LDL^T$ over $\mathbb{Q}$, $n \le 7$ (1361 satisfying \cref{eq:ldlt-crit}) & 2000 & 0 \\
    Canonical form $=$ rank-jump profile over $\mathbb{Q}$ (\Cref{prop:canonical}) & 500 & 0 \\
    Exhaustive LU over $GF(2)$, $n \le 3$, vs.\ brute force & 530 & 0 \\
    Random LU over $GF(2)$, $n = 4, \ldots, 6$ & 5000 & 0 \\
    Exhaustive $LDL^T$ over $GF(2)$, symmetric, $n \le 4$, vs.\ brute force & 1098 & 0 \\
    Exhaustive LU over $GF(3)$: $n = 2$ vs.\ brute force; all of $n = 3$ & 19764 & 0 \\
    Exhaustive $LDL^T$ over $GF(3)$, symmetric, $n = 3$, vs.\ brute force & 729 & 0 \\
    \midrule
    Total & 31639 & 0 \\
    \bottomrule
  \end{tabular}
  \vspace{5pt}
  \caption{Exact-arithmetic validation results.  "vs.\ brute force" means
    the criterion was compared, on every matrix of the class, against
    exhaustive enumeration of all triangular factorizations --- an
    independent ground truth.  Every satisfying case also verifies the
    constructed factors exactly.}\label{tab:validation}
\end{table}

The characteristic-$2$ rows deserve emphasis: they exercise precisely the
regime in which signed or real-specific arguments would fail (cf.\
\Cref{rem:signed}), and confirm the field-general constructions,
including the $\beta$-weighted pair blocks and the corner-free thread
recursion of \Cref{lem:telescope}.

\section{Conclusion}\label{sec:conclusion}

We gave a unified treatment of the existence of triangular factorizations
without pivoting for possibly singular matrices over an arbitrary field.
The unsymmetric criterion (\Cref{thm:lu}) recovers the theorem of Okunev
and Johnson; the symmetric criterion (\Cref{thm:ldlt}) --- existence of
$A = LDL^T$ with truly diagonal $D$ and possibly singular $L$ --- appears
to be new, as does the corollary that for symmetric matrices unpivoted
$LU$ and $LDL^T$ existence coincide.  The proof method is uniform: reduce
to a matching normal form by triangular equivalence or congruence,
transfer the nullity criterion to the normal form where it becomes a
prefix condition on the matching, and factor the normal form explicitly
--- by greedy slot assignment in the unsymmetric case, by telescoping
threads rooted at isolated zero indices in the symmetric case.  The
normal forms tie the existence theory to the generalized Bruhat
decomposition and rank profile matrix, and, in the symmetric case, to the
factorization of Van Dooren and Strang and the congruence canonical forms
of Szechtman.  Everything is constructive, and the resulting $O(n^3)$
algorithm was validated exhaustively in exact arithmetic.

Several directions remain open.
\begin{itemize}
  \item \emph{Numerical stability.}  The construction selects pivots for
    algebraic, not numerical, reasons; a backward-stable variant over
    $\mathbb{R}$ (perhaps with size-aware choices among the admissible
    threads and slots, which are far from unique) is an open problem, as
    in~\cite{darve_lu_2026}.
  \item \emph{Multiplicity.}  Dopico, Johnson, and
    Molera~\cite{dopico_multiple_2006} parametrized all LU factorizations
    of a singular matrix with $L$ unit lower triangular.  Parametrizing
    all unpivoted $LDL^T$ factorizations --- equivalently, all thread
    systems and weight choices --- would quantify the freedom our greedy
    constructions exploit.
  \item \emph{Structure of the factors.}  The rank-revealing and sparsity
    results of~\cite{darve_lu_2026} for LU should have $LDL^T$ analogues;
    the thread construction suggests that the columns of $L$ can be
    localized to their threads.
  \item \emph{Hermitian and other variants.}  Over $\mathbb{C}$ with
    conjugate transposition, congruence becomes $\ast$-congruence and the
    pair blocks acquire Hermitian weights; we expect the theory to carry
    over with $\beta$ real, but have not pursued it.  Rectangular and
    block versions of the criteria are also natural.
\end{itemize}


\appendix
\section{Reference Python implementation}\label{sec:appendix-code}

\Cref{lst:matching} lists the complete implementation used in
\Cref{sec:validation}: the field abstraction (rationals via
\texttt{fractions.Fraction}, and prime fields $GF(p)$), exact rank and the
two nullity criteria, the two reductions to matching normal forms
(\Cref{lem:reduction-lu,lem:reduction-sym}), the greedy slot assignment
(\Cref{thm:lu-matching}), and the thread construction
(\Cref{lem:threads,lem:telescope}).  The entry points
\texttt{lu\_factorization} and \texttt{ldlt\_factorization} return exact
triangular factors, or \texttt{None} exactly when the corresponding
criterion fails.  The test driver \texttt{validate.py} (not listed;
available with the source of this paper) implements the four validation
layers of \Cref{sec:validation}, including the brute-force enumerations
over $GF(2)$ and $GF(3)$.

\lstinputlisting[style=PythonStyle,
  caption={Constructive unpivoted LU and $LDL^T$ factorization via
    matching normal forms, in exact field-generic arithmetic.},
  label={lst:matching},
  rangeprefix=\#\ ,
  linerange=scifig_start-scifig_end,
  includerangemarker=false]{code/matching_factorization.py}


\clearpage
\phantomsection
\bibliographystyle{siamplain}
\bibliography{ref}

\end{document}
